% =====================================================================
%  Tesi di Laurea Magistrale — Engineering in Computer Science
%  Sapienza Università di Roma
%
%  Nuova stesura della tesi del 2022 (legacy/tesi-magistrale/), basata sulla
%  reimplementazione del repository (src/rsseg/, docs/).
%
%  Compilazione:  latexmk            (pdflatex + bibtex, vedi latexmkrc)
%  Pulizia:       latexmk -c
% =====================================================================

% Stessa classe della tesi del 2022 (A4 fronte-retro; il corpo 11pt è il default su A4).
\documentclass[a4paper,twoside]{sapthesis}

% --- Modalità bozza: mostra note \todo, segnaposto e indice delle note -------
\newif\ifbozza
\bozzatrue        % \bozzafalse per la versione da consegnare

\newcommand{\TitoloTesi}{Segmentazione semantica di immagini telerilevate ad altissima risoluzione}

% =====================================================================
%  Pacchetti
%  La classe Sapthesis carica già: geometry, fontenc (T1), lmodern,
%  caption (etichetta piccola e in grassetto, separatore "."), graphicx,
%  color, booktabs, amsmath, etoolbox, fancyhdr.
% =====================================================================

% --- Lingua e testo ----------------------------------------------------
\usepackage[italian]{babel}
\usepackage[autostyle,italian=guillemets]{csquotes}
\usepackage{microtype}
\usepackage{icomma}        % virgola decimale in modo matematico: $0,604$
\usepackage{enumitem}
\setlist{itemsep=0.15em, parsep=0.15em, topsep=0.3em}

% Paragrafi separati da una riga vuota e senza rientro, come nella tesi del 2022.
\usepackage[skip=0.75\baselineskip plus 2pt, indent=0pt]{parskip}

% --- Colori (xcolor sostituisce color, caricato dalla classe) ------------
\usepackage[table,dvipsnames]{xcolor}
% =====================================================================
%  Palette colori
%  Ricavata dal PDF della tesi del 2022 (legacy/tesi-magistrale/), leggendo
%  i valori nei content stream e nelle annotazioni dei link:
%    - rosso Sapienza RGB(130,36,51): frontespizio (definito dalla classe
%      Sapthesis come `sapred', ridefinito qui per usarlo già nel preambolo);
%    - cornici dei link di hyperref: rosso [1 0 0] per i riferimenti interni,
%      verde [0 1 0] per le citazioni, ciano [0 1 1] per gli URL.
%  Tutti gli altri colori sono tinte di questi, oppure i colori ufficiali
%  delle classi ISPRS, che servono a leggere maschere e legende.
% =====================================================================

% --- Colori di identità (identici alla tesi del 2022) -----------------
\definecolor{SapRosso}{rgb}{0.5098039,0.1411765,0.2}   % = sapred, #822433
\definecolor{LinkRosso}{rgb}{1,0,0}                    % riferimenti interni
\definecolor{CitaVerde}{rgb}{0,1,0}                    % citazioni
\definecolor{UrlCiano}{rgb}{0,1,1}                     % URL

% --- Tinte di supporto per schemi e tabelle ---------------------------
\colorlet{SapRossoScuro}{SapRosso!80!black}
\colorlet{SapRossoMedio}{SapRosso!60}
\colorlet{SapRossoTenue}{SapRosso!30}
\colorlet{SapRossoChiaro}{SapRosso!12}
\colorlet{SapRossoVelo}{SapRosso!5}
\colorlet{SapGrigio}{black!55}
\colorlet{SapGrigioTenue}{black!25}
\colorlet{SapGrigioChiaro}{black!8}

% --- Classi ISPRS (palette ufficiale delle maschere di Potsdam/Vaihingen)
\definecolor{ClsImpervious}{RGB}{255,255,255}   % superfici impermeabili
\definecolor{ClsBuilding}{RGB}{0,0,255}         % edifici
\definecolor{ClsLowVeg}{RGB}{0,255,255}         % vegetazione bassa
\definecolor{ClsTree}{RGB}{0,255,0}             % alberi
\definecolor{ClsCar}{RGB}{255,255,0}            % auto
\definecolor{ClsClutter}{RGB}{255,0,0}          % clutter / sfondo

% --- Ruoli semantici negli schemi --------------------------------------
\colorlet{ColEncoder}{SapRossoChiaro}     % encoder / backbone
\colorlet{ColDecoder}{SapGrigioChiaro}    % decoder
\colorlet{ColAttn}{SapRossoTenue}         % blocchi di attenzione
\colorlet{ColQuant}{SapRosso}             % quantizzazione vettoriale
\colorlet{ColFreccia}{SapGrigio}          % frecce di flusso
\colorlet{ColGradiente}{SapRosso}         % percorsi del gradiente

% --- Split dei dati -----------------------------------------------------
\colorlet{SplitTrain}{SapRossoMedio}
\colorlet{SplitVal}{SapRossoTenue}
\colorlet{SplitTest}{SapGrigio}
\colorlet{SplitEscl}{white}

% --- Righe evidenziate e segnaposto -------------------------------------
\colorlet{RigaRiferimento}{SapRossoVelo}  % righe di riferimento nelle tabelle
\colorlet{ColTodo}{SapRossoChiaro}        % sfondo delle note \todo
\colorlet{ColTodoBordo}{SapRosso}


% --- Matematica ----------------------------------------------------------
\usepackage{amssymb}
\usepackage{amsthm}
\usepackage{mathtools}
\usepackage{bm}

% --- Tabelle -------------------------------------------------------------
\usepackage{array}
\usepackage{multirow}
\usepackage{makecell}
\usepackage{tabularx}
\usepackage{siunitx}
\sisetup{output-decimal-marker={,}, group-separator={.}, group-minimum-digits=4,
         detect-weight=true, detect-family=true}
\renewcommand{\arraystretch}{1.15}

% --- Figure e schemi -----------------------------------------------------
\usepackage{float}          % [H]: posizione fissa (usata nei capitoli con segnaposto)
\usepackage{subcaption}
\usepackage{tikz}
\usetikzlibrary{positioning, arrows.meta, calc, fit, backgrounds,
                shapes.geometric, shapes.misc, matrix, decorations.pathreplacing,
                patterns, shadows}
\usepackage{pgfplots}
\pgfplotsset{compat=1.18}

% --- Algoritmi -----------------------------------------------------------
\usepackage{algorithm}
\usepackage{algpseudocode}
\floatname{algorithm}{Algoritmo}
\renewcommand{\listalgorithmname}{Elenco degli algoritmi}
\algrenewcommand\algorithmicrequire{\textbf{Input:}}
\algrenewcommand\algorithmicensure{\textbf{Output:}}
\algrenewcommand\algorithmicfor{\textbf{per}}
\algrenewcommand\algorithmicdo{\textbf{esegui}}
\algrenewcommand\algorithmicend{\textbf{fine}}
\algrenewcommand\algorithmicif{\textbf{se}}
\algrenewcommand\algorithmicthen{\textbf{allora}}
\algrenewcommand\algorithmicreturn{\textbf{restituisci}}
\algnewcommand\algorithmicforeach{\textbf{per ogni}}
\algdef{S}[FOR]{ForEach}[1]{\algorithmicforeach\ #1\ \algorithmicdo}

% --- Note di lavoro e segnaposto -----------------------------------------
% Per la versione finale: \bozzafalse in main.tex (le note spariscono).
\ifbozza
  \usepackage[colorinlistoftodos, italian, textsize=footnotesize,
              color=ColTodo, bordercolor=ColTodoBordo, linecolor=ColTodoBordo]{todonotes}
\else
  \usepackage[disable]{todonotes}
\fi
\setlength{\marginparwidth}{2.2cm}

% --- Collegamenti (stesse cornici colorate della tesi del 2022) -----------
\usepackage{hyperref}
\hypersetup{
  pdftitle    = {\TitoloTesi},
  pdfauthor   = {Matteo Campana},
  pdfsubject  = {Tesi di Laurea Magistrale in Engineering in Computer Science},
  pdfkeywords = {segmentazione semantica, telerilevamento, attenzione lineare,
                 quantizzazione vettoriale, apprendimento auto-supervisionato, DINO},
  colorlinks      = false,
  pdfborder       = {0 0 1},
  linkbordercolor = LinkRosso,
  citebordercolor = CitaVerde,
  urlbordercolor  = UrlCiano,
  bookmarksnumbered = true,
}
\usepackage[italian, noabbrev]{cleveref}
\crefname{algorithm}{algoritmo}{algoritmi}
\Crefname{algorithm}{Algoritmo}{Algoritmi}
\crefname{ipotesi}{ipotesi}{ipotesi}
\Crefname{ipotesi}{Ipotesi}{Ipotesi}

% =====================================================================
%  Comandi e ambienti
% =====================================================================

% --- Enunciati -------------------------------------------------------------
\theoremstyle{definition}
\newtheorem{definizione}{Definizione}[chapter]
\newtheorem{ipotesi}{Ipotesi}
\theoremstyle{plain}
\newtheorem{proposizione}[definizione]{Proposizione}
\theoremstyle{remark}
\newtheorem{osservazione}[definizione]{Osservazione}
\renewcommand{\qedsymbol}{$\blacksquare$}
\crefname{definizione}{definizione}{definizioni}
\Crefname{definizione}{Definizione}{Definizioni}
\crefname{proposizione}{proposizione}{proposizioni}
\Crefname{proposizione}{Proposizione}{Proposizioni}
\crefname{osservazione}{osservazione}{osservazioni}
\Crefname{osservazione}{Osservazione}{Osservazioni}

% --- Notazione matematica ----------------------------------------------------
\newcommand{\R}{\mathbb{R}}
\newcommand{\N}{\mathbb{N}}
\newcommand{\E}{\mathbb{E}}
\newcommand{\Ind}{\mathbf{1}}
\DeclareMathOperator*{\argmin}{arg\,min}
\DeclareMathOperator*{\argmax}{arg\,max}
\DeclareMathOperator{\softmax}{softmax}
\DeclareMathOperator{\sg}{sg}                 % stop-gradient
\DeclareMathOperator{\GAP}{GAP}               % global average pooling
\DeclareMathOperator{\PPL}{PPL}               % perplexity del codebook
\DeclareMathOperator{\KL}{D_{KL}}
\newcommand{\Hent}{\mathrm{H}}
\newcommand{\norm}[1]{\left\lVert #1 \right\rVert}
\newcommand{\abs}[1]{\left\lvert #1 \right\rvert}
\newcommand{\T}{^{\top}}
\newcommand{\bigO}[1]{\mathcal{O}\!\left(#1\right)}
\newcommand{\vect}[1]{\bm{#1}}
\newcommand{\hatq}{\hat{\vect{q}}}
\newcommand{\hatk}{\hat{\vect{k}}}

% --- Nomi ricorrenti -----------------------------------------------------------
\newcommand{\maresunet}{MAResU-Net}
\newcommand{\maresunetvq}{MAResU-Net-VQ}
\newcommand{\rsseg}{\texttt{rsseg}}
\newcommand{\RQ}[1]{\textbf{RQ#1}}
\newcommand{\codice}[1]{\texttt{#1}}
\newcommand{\tile}[1]{\texttt{#1}}

% --- Quadratino con il colore di una classe ISPRS (per le legende) -------------
\newcommand{\ClasseBox}[1]{%
  \tikz[baseline=-0.6ex]\draw[black!60, line width=0.3pt, fill=#1] (0,0) rectangle (0.8em,0.8em) ++(0,-0.4em);}
\newcommand{\ClasseBoxX}[1]{%
  \tikz[baseline=0.05em]\draw[black!60, line width=0.3pt, fill=#1] (0,0) rectangle (0.8em,0.8em);}

% --- Segnaposto per le parti sperimentali ------------------------------------
% \Segnaposto{...}  nota in linea, ben visibile, per contenuti ancora da produrre
% \ND               valore non ancora disponibile in una cella di tabella
% \FiguraMancante{descrizione}{altezza}  riquadro al posto di una figura
\newcommand{\Segnaposto}[1]{\todo[inline]{#1}}
\newcommand{\ND}{\textcolor{SapRosso}{\textit{n.d.}}}
\newcommand{\FiguraMancante}[2][4cm]{%
  \ifbozza
    \begin{tikzpicture}
      \draw[SapRosso, dashed, line width=0.8pt, fill=SapRossoVelo]
        (0,0) rectangle (0.95\linewidth, #1);
      \node[text width=0.8\linewidth, align=center, font=\small\itshape, text=SapRossoScuro]
        at (0.475\linewidth, #1/2) {Figura da inserire\\[0.3em]#2};
    \end{tikzpicture}%
  \else
    \fbox{\parbox[c][#1][c]{0.95\linewidth}{\centering\small\itshape Figura da inserire: #2}}%
  \fi}

% --- Stili TikZ condivisi dagli schemi -------------------------------------
\tikzset{
  >={Stealth[length=2.2mm]},
  blocco/.style={draw=SapRossoScuro, line width=0.5pt, rounded corners=1.5pt,
                 align=center, font=\footnotesize, minimum height=7mm, inner sep=2.5pt},
  enc/.style={blocco, fill=ColEncoder},
  dec/.style={blocco, fill=ColDecoder, draw=SapGrigio},
  attn/.style={blocco, fill=ColAttn},
  quant/.style={blocco, fill=ColQuant, text=white, draw=SapRossoScuro},
  dato/.style={blocco, draw=SapGrigio, fill=white, rounded corners=0pt},
  nota/.style={font=\scriptsize\itshape, text=SapGrigio, align=center},
  flusso/.style={->, draw=ColFreccia, line width=0.6pt},
  grad/.style={->, draw=ColGradiente, line width=0.6pt, dashed},
  skipc/.style={->, draw=SapRosso, line width=0.6pt},
  gruppo/.style={draw=SapGrigioTenue, dashed, rounded corners=3pt, inner sep=6pt},
}


% --- Frontespizio (dati come nella tesi del 2022) ------------------------
\title{\TitoloTesi}
\subtitle{Attenzione lineare, quantizzazione vettoriale e auto-distillazione:\\
          una rivalutazione riproducibile}
\author{Matteo Campana}
\IDnumber{1692599}
\courseorganizer{Facolt\`a di Ingegneria dell'Informazione, Informatica e Statistica}
\course{Corso di Laurea Magistrale in Engineering in Computer Science -- Ingegneria Informatica}
\AcademicYear{2026/2027}   % TODO: verificare l'anno accademico della discussione
\advisor{Prof. Aurelio Uncini}
\coadvisor{Dr. Simone Scardapane}
\coadvisor{Dr. Valerio Marsocci}
\authoremail{campana.1692599@studenti.uniroma1.it}
\copyyear{2026}

\begin{document}

\frontmatter
\maketitle

% --- Dedica -------------------------------------------------------------------
\cleardoublepage
\thispagestyle{empty}
\vspace*{0.25\textheight}
\begin{flushright}
  \itshape
  Dedicato a tutte le persone che, nonostante gli ostacoli,\\
  mi hanno impedito di lasciare questo percorso.
\end{flushright}

% =====================================================================
%  Abstract
% =====================================================================
% La classe, in italiano, intitola l'ambiente "Sommario": qui si usa "Abstract".
\begin{abstract}[Abstract]
La segmentazione semantica di ortofoto aeree ad altissima risoluzione assegna una classe di
copertura del suolo (superfici impermeabili, edifici, vegetazione bassa, alberi, auto, clutter) a
ogni pixel dell'immagine. Il compito è reso difficile da un'elevata variabilità intra-classe
unita a una bassa variabilità inter-classe, dalle ombre e dalle occlusioni e, soprattutto, dal
costo delle annotazioni, che restano poche e richiedono personale esperto.

Una prima tesi, discussa nel 2022, ha studiato due estensioni dell'architettura \maresunet{}
(encoder ResNet, decoder U-Net e \emph{Linear Attention Mechanism} nelle connessioni di salto):
un collo di bottiglia quantizzato con un codebook di tipo VQ-VAE e un pre-addestramento
auto-supervisionato dell'encoder per auto-distillazione (DINO) sul dataset ISPRS Potsdam.
Entrambe peggioravano la baseline pre-addestrata su ImageNet. Una revisione sistematica di quel
lavoro mostra però che il risultato non è conclusivo: la classe clutter era esclusa dalla
funzione di perdita, le metriche erano mediate per batch, la validation probabilmente non era
indipendente dal training e il pre-addestramento si discostava in punti sostanziali dal metodo
originale.

Questa tesi ripropone le stesse domande con un'implementazione nuova, verificata e
riproducibile, e con un protocollo confrontabile con la letteratura: split standard per tile
intere, valutazione su tile complete con etichette piene ed \emph{eroded}, matrice di confusione
globale e più seed per configurazione. Le domande di ricerca riguardano il contributo
dell'attenzione lineare (\RQ{1}), l'effetto di un codebook monitorato e protetto dal collasso
(\RQ{2}), il valore di un pre-addestramento DINO fedele all'originale al variare della quantità
di etichette (\RQ{3}), l'adeguatezza delle trasformazioni di \emph{data augmentation} alle immagini
aeree (\RQ{4}), il compromesso tra accuratezza e costo computazionale su una singola GPU da
\SI{8}{\giga\byte} (\RQ{5}) e l'uso dell'altezza normalizzata sul terreno (nDSM) come quarto
canale di ingresso (\RQ{6}). Il secondo benchmark è il dataset ISPRS Vaihingen.

Il lavoro sviluppa il quadro teorico dei tre meccanismi: la derivazione dell'attenzione a costo
lineare nel numero di pixel, la quantizzazione vettoriale come K-means online con stimatore
\emph{straight-through}, l'analisi del collasso nell'auto-distillazione. Definisce inoltre un
disegno sperimentale che collega ogni domanda di ricerca a una campagna di esperimenti
eseguibile e riprendibile.

\Segnaposto{Aggiungere una sintesi quantitativa dei risultati (mIoU sul test di Potsdam e
Vaihingen, effetto di VQ, DINO e nDSM) quando le campagne saranno completate.}
\end{abstract}

% =====================================================================
%  Ringraziamenti
% =====================================================================
\begin{acknowledgments}
Ringrazio il Prof.\ Aurelio Uncini, il Dr.\ Simone Scardapane e il Dr.\ Valerio Marsocci per la
guida, i commenti sulla configurazione sperimentale e i suggerimenti durante la preparazione
della tesi.

Ringrazio l'International Society for Photogrammetry and Remote Sensing (ISPRS) per aver reso
disponibili i dataset Potsdam, Vaihingen e Toronto.

\Segnaposto{Completare i ringraziamenti personali.}
\end{acknowledgments}


\tableofcontents
\listoffigures
\listoftables
\ifbozza
  \cleardoublepage
  \listoftodos[Elenco delle parti da completare]
\fi

\mainmatter
% =====================================================================
%  Capitolo 1 — Introduzione
% =====================================================================
\chapter{Introduzione}
\label{cap:introduzione}

\section{Contesto}
\label{sec:contesto}

Le immagini telerilevate ad altissima risoluzione (\emph{Very High Resolution}, VHR), riprese
da aereo o da satellite con una distanza di campionamento al suolo (\emph{Ground Sampling
Distance}, GSD) di pochi centimetri, sono oggi una fonte primaria di informazione per il
monitoraggio delle aree urbane, la pianificazione dell'uso del suolo, la tutela ambientale e la
gestione del traffico~\cite{zhu2017dlrs,cheng2016survey}. Per trasformare queste immagini in
informazione utilizzabile occorre sapere \emph{che cosa} c'è in ogni punto della scena: è il
compito della \emph{segmentazione semantica}, che assegna a ciascun pixel una classe scelta da un
insieme fissato a priori.

Rispetto alla classificazione di un'immagine intera o al rilevamento di oggetti con riquadri, la
segmentazione semantica è una predizione \emph{densa}: per un'ortofoto di $6000\times6000$ pixel
il modello produce $3{,}6\cdot10^{7}$ decisioni, che devono rispettare i contorni degli oggetti e
rimanere coerenti su scala urbana. Dal 2015 il problema è affrontato con reti neurali
convoluzionali completamente convoluzionali~\cite{long2015fcn} e con architetture
encoder-decoder come U-Net~\cite{ronneberger2015unet}. Negli anni successivi queste architetture
sono state arricchite con convoluzioni dilatate~\cite{chen2018deeplab}, meccanismi di
attenzione~\cite{wang2018nonlocal,fu2019danet} e, più di recente, con encoder di tipo
transformer~\cite{dosovitskiy2021vit,liu2021swin,xie2021segformer}.

Il banco di prova di riferimento per la segmentazione di scene urbane aeree è il \emph{2D
Semantic Labeling Contest} dell'ISPRS~\cite{rottensteiner2012isprs,isprs2dlabeling}, con i
dataset \textbf{Potsdam} (ortofoto RGB a \SI{5}{\centi\meter}) e \textbf{Vaihingen} (falso colore
IR-R-G a \SI{9}{\centi\meter}), annotati in sei classi: superfici impermeabili
\ClasseBox{ClsImpervious}, edifici \ClasseBox{ClsBuilding}, vegetazione bassa
\ClasseBox{ClsLowVeg}, alberi \ClasseBox{ClsTree}, auto \ClasseBox{ClsCar} e clutter
\ClasseBox{ClsClutter}.

\section{Le difficoltà del dominio}
\label{sec:difficolta}

Il telerilevamento urbano VHR combina quattro difficoltà, che motivano le scelte di questa tesi
(\cref{fig:filo-logico}).

\begin{description}[style=nextline, leftmargin=1.2em]
  \item[Variabilità intra-classe alta e inter-classe bassa.] Gli edifici hanno materiali,
  forme e colori molto diversi tra loro, mentre tetti, marciapiedi e strade in cemento possono
  essere quasi indistinguibili nel solo spettro visibile. Lo stesso vale per alberi e
  vegetazione bassa, che si distinguono soprattutto per l'altezza.
  \item[Condizioni di ripresa.] Ombre degli edifici alti, foschia e differenze di
  illuminazione alterano l'aspetto delle superfici e introducono confusioni sistematiche.
  \item[Scarsità di etichette.] A differenza dei grandi dataset di immagini naturali, le
  annotazioni pixel per pixel di immagini aeree sono poche, costose e richiedono esperti: due
  delle collezioni ISPRS più usate contano 38 (Potsdam) e 33 (Vaihingen) aree annotate.
  \item[Contesto ampio.] Classificare un pixel richiede di guardare l'oggetto e il suo intorno:
  un'auto si riconosce dalla strada su cui si trova, un cortile dagli edifici che lo
  circondano. Aggregare contesto su immagini grandi è però costoso. L'attenzione
  \emph{dot-product} ha un costo quadratico nel numero di posizioni, già proibitivo su un
  singolo ritaglio di $512\times512$ pixel (\cref{sec:attenzione-dot}).
\end{description}

\begin{figure}[t]
  \centering
  \resizebox{\linewidth}{!}{% Schema: dai problemi del dominio ai meccanismi studiati e alle domande di ricerca.
\begin{tikzpicture}[
    node distance=4mm and 9mm,
    prob/.style={dato, text width=3.5cm, minimum height=10mm},
    mecc/.style={attn, text width=3.4cm, minimum height=10mm},
    domanda/.style={blocco, fill=white, draw=SapRosso, text width=2.6cm, minimum height=10mm},
  ]
  % Colonna 1: problemi
  \node[prob] (p1) {Immagini molto grandi\\{\scriptsize $6000\times6000$ px, contesto ampio}};
  \node[prob, below=of p1] (p2) {Variabilit\`a intra-classe alta,\\inter-classe bassa};
  \node[prob, below=of p2] (p3) {Etichette scarse e costose};
  \node[prob, below=of p3] (p4) {Ambiguit\`a spettrale\\{\scriptsize tetto vs strada, albero vs prato}};

  % Colonna 2: meccanismi
  \node[mecc, right=of p1] (m1) {Attenzione lineare\\{\scriptsize LAM, costo $\mathcal{O}(N)$}};
  \node[mecc, right=of p2] (m2) {Bottleneck discreto\\{\scriptsize codebook VQ}};
  \node[mecc, right=of p3] (m3) {Auto-distillazione\\{\scriptsize DINO su dati non etichettati}};
  \node[mecc, right=of p4] (m4) {Altezza sul terreno\\{\scriptsize nDSM come 4\textordmasculine{} canale}};

  % Colonna 3: domande di ricerca
  \node[domanda, right=of m1] (r1) {\RQ{1}, \RQ{5}\\{\scriptsize contributo e costo}};
  \node[domanda, right=of m2] (r2) {\RQ{2}\\{\scriptsize codebook sano?}};
  \node[domanda, right=of m3] (r3) {\RQ{3}, \RQ{4}\\{\scriptsize poche etichette, augmentation}};
  \node[domanda, right=of m4] (r4) {\RQ{6}\\{\scriptsize fusione precoce}};

  \foreach \i in {1,...,4}{
    \draw[flusso] (p\i) -- (m\i);
    \draw[flusso] (m\i) -- (r\i);
  }
  % legami secondari
  \draw[flusso, densely dotted] (p2.east) -- (m4.north west);
  \draw[flusso, densely dotted] (p3.east) -- (m2.south west);

  % Intestazioni
  \node[nota, above=2mm of p1, font=\footnotesize\bfseries, text=SapRossoScuro] {Problema};
  \node[nota, above=2mm of m1, font=\footnotesize\bfseries, text=SapRossoScuro] {Meccanismo};
  \node[nota, above=2mm of r1, font=\footnotesize\bfseries, text=SapRossoScuro] {Domanda};

  % Riquadro: estensioni rispetto al 2022
  \begin{scope}[on background layer]
    \node[gruppo, fit=(m4)(r4), label={[nota]below:estensione rispetto alla tesi del 2022}] {};
  \end{scope}
\end{tikzpicture}
}
  \caption[Dai problemi del dominio alle domande di ricerca]{Dai problemi del dominio ai
    meccanismi studiati e alle domande di ricerca. Le linee punteggiate indicano legami
    secondari: l'altezza aiuta a separare classi spettralmente simili e la quantizzazione nasce
    dall'idea di \emph{visual words} usata dal pre-addestramento OBoW.}
  \label{fig:filo-logico}
\end{figure}

\section{Dalla tesi del 2022 a questo lavoro}
\label{sec:da-2022}

Questa tesi riprende il lavoro svolto nel 2022~\cite{campana2022tesi}, a sua volta continuazione
di MARE~\cite{marsocci2021mare}. MARE ha mostrato che l'architettura
\maresunet{}~\cite{li2022maresunet}, cioè un encoder ResNet~\cite{he2016resnet} e un decoder
U-Net con un \emph{Linear Attention Mechanism}~(LAM)~\cite{li2020lam} nelle connessioni di
salto, beneficia sul dataset Vaihingen di un pre-addestramento auto-supervisionato basato su
\emph{bag of visual words}~(OBoW)~\cite{gidaris2021obow}. La tesi del 2022 ha esteso
quell'idea in tre direzioni: un dataset più grande (Potsdam), un metodo di auto-distillazione
più recente (DINO~\cite{caron2021dino}) e un collo di bottiglia esplicitamente discreto, con un
codebook di quantizzazione vettoriale~\cite{vandenoord2017vqvae} ispirato al vocabolario di
\emph{visual words} di OBoW.

Il risultato sperimentale di allora è stato negativo per entrambe le estensioni: la
configurazione migliore restava la \maresunet{} con encoder ResNet-18 pre-addestrato su
ImageNet (mIoU di test $\approx 0{,}604$). Una revisione sistematica del protocollo e del codice
originale (\cref{sec:lettura-critica}) mostra tuttavia che quel risultato \emph{non è
conclusivo}. La classe clutter non riceveva gradiente, le metriche erano mediate per batch, la
validation probabilmente non era indipendente dal training e il pre-addestramento si
discostava da DINO in punti sostanziali. Le domande poste nel 2022 restano quindi aperte e
meritano una risposta basata su un protocollo corretto.

\section{Obiettivi}
\label{sec:obiettivi}

L'obiettivo della tesi è \textbf{rivalutare in modo riproducibile e confrontabile con la
letteratura} il contributo di tre meccanismi (attenzione lineare, quantizzazione vettoriale e
auto-distillazione) alla segmentazione semantica di immagini aeree urbane, ed estendere
l'analisi all'uso dell'informazione altimetrica. In dettaglio:

\begin{enumerate}
  \item formalizzare i meccanismi studiati e derivarne le proprietà rilevanti per il dominio:
  costo computazionale, capacità informativa del collo di bottiglia, modi di collasso
  dell'auto-distillazione (\cref{cap:teoria});
  \item definire un protocollo sperimentale che elimini i fattori confondenti individuati nel
  lavoro del 2022 (\cref{cap:metodologia});
  \item rispondere con più seed e intervalli di variabilità alle domande di ricerca della
  \cref{sec:domande} (\cref{cap:esperimenti});
  \item quantificare il compromesso tra accuratezza e costo su hardware consumer (una GPU da
  \SI{8}{\giga\byte}), in linea con la motivazione energetica già presente nel 2022.
\end{enumerate}

\section{Domande di ricerca}
\label{sec:domande}

\begin{description}[style=nextline, leftmargin=1.2em]
  \item[\RQ{1} -- Attenzione lineare.] Su Potsdam, con un protocollo standard, quanto rende
  \maresunet{} rispetto a una rete con la stessa struttura ma connessioni di salto senza
  attenzione, a parità di encoder?
  \item[\RQ{2} -- Quantizzazione vettoriale.] Un collo di bottiglia quantizzato, con un codebook
  \emph{sano} (perplexity alta, pochi codici inutilizzati), migliora la segmentazione? Conta dove
  la quantizzazione viene applicata: solo al collo di bottiglia o anche alle connessioni di
  salto?
  \item[\RQ{3} -- Auto-distillazione e poche etichette.] Un pre-addestramento DINO fedele
  all'originale, su immagini di dominio non etichettate, eguaglia o supera ImageNet? Come cambia
  il confronto quando si usa solo l'1\%, il 10\% o il 25\% delle etichette?
  \item[\RQ{4} -- Augmentation per immagini aeree.] Le trasformazioni cromatiche
  \enquote{forti} di DINO sono adatte alle immagini aeree, o conviene una variante attenuata e
  invariante alle rotazioni?
  \item[\RQ{5} -- Efficienza.] Qual è il compromesso tra accuratezza e costo (parametri,
  memoria, throughput) tra encoder ResNet-18, ResNet-34 e ResNet-50 e tra encoder
  convoluzionali e transformer?
  \item[\RQ{6} -- Altezza sul terreno.] L'nDSM come quarto canale di ingresso migliora la
  segmentazione, in particolare per le coppie edifici/superfici impermeabili e alberi/vegetazione
  bassa? Il risultato cambia se l'nDSM è quello ufficiale o quello stimato dal DSM?
\end{description}

Le domande da \RQ{1} a \RQ{5} riformulano le ipotesi del 2022. La \RQ{6} e il secondo benchmark
(Vaihingen) sono estensioni introdotte in questo lavoro.

\section{Contributi}
\label{sec:contributi}

\begin{itemize}
  \item Una \textbf{lettura critica} della tesi del 2022 che separa le conclusioni giustificate
  dai dati da quelle dovute a difetti di protocollo o di implementazione (\cref{cap:premesse}).
  \item Una \textbf{trattazione teorica unificata} dei meccanismi studiati. Comprende la
  derivazione del LAM come attenzione a nucleo con costo $\bigO{N}$ e una sua proprietà di
  \enquote{basso contrasto}, l'interpretazione dell'aggiornamento EMA del codebook come K-means
  online, l'analisi dell'aggiramento del collo di bottiglia da parte delle connessioni di salto,
  i modi di collasso dell'auto-distillazione e l'effetto di una classe esclusa dalla perdita
  (\cref{cap:teoria}).
  \item Un \textbf{protocollo sperimentale} con split standard per tile intere, valutazione su
  tile complete con etichette piene ed \emph{eroded}, matrice di confusione globale, più seed e
  un regime a poche etichette con budget di iterazioni costante (\cref{cap:metodologia}).
  \item Un'\textbf{implementazione nuova e verificata} (pacchetto \rsseg{}, Keras~3 su backend
  PyTorch, tracciamento con MLflow). Comprende campagne di esperimenti riprendibili, 109 test
  automatici e la verifica di equivalenza numerica con i modelli di riferimento
  (\cref{sec:implementazione}).
  \item I \textbf{risultati sperimentali} che rispondono alle domande di ricerca
  (\cref{cap:esperimenti}).
  \Segnaposto{Riassumere i risultati principali in una riga per domanda di ricerca, quando
  disponibili.}
\end{itemize}

\section{Struttura della tesi}
\label{sec:struttura}

Il \cref{cap:premesse} inquadra il lavoro nello stato dell'arte e analizza criticamente la tesi
del 2022. Il \cref{cap:teoria} sviluppa il quadro teorico: formulazione del problema, encoder-decoder
residuali, attenzione lineare, quantizzazione vettoriale, auto-distillazione, fusione
dell'nDSM, metriche e inferenza su tile intere. Il \cref{cap:metodologia} descrive dati,
protocolli, configurazioni di addestramento, implementazione e analisi statistica. Il
\cref{cap:esperimenti} è dedicato agli esperimenti, organizzati per domanda di ricerca. Il
\cref{cap:prospettive} discute lo stato del lavoro, le attività rimanenti e le prospettive di
ricerca.

% =====================================================================
%  Capitolo 2 — Premesse: stato dell'arte e lavoro del 2022
% =====================================================================
\chapter{Premesse}
\label{cap:premesse}

Questo capitolo colloca il lavoro nel contesto della letteratura e ne fissa il punto di
partenza. Le prime sezioni ripercorrono le quattro linee di ricerca su cui la tesi si appoggia:
segmentazione semantica, attenzione efficiente, rappresentazioni discrete e apprendimento
auto-supervisionato. Le ultime riassumono la tesi del 2022 e ne propongono una lettura critica,
da cui discendono le scelte di metodo dei capitoli successivi.

\section{Segmentazione semantica con reti profonde}
\label{sec:stato-segmentazione}

Le \emph{Fully Convolutional Networks} (FCN) di Long et al.~\cite{long2015fcn} hanno mostrato
che una rete di classificazione può produrre una predizione densa sostituendo gli strati
completamente connessi con convoluzioni e ricampionando le mappe di attivazione alla risoluzione
dell'ingresso. Il limite principale di questo schema è la perdita di dettaglio spaziale dovuta
al sottocampionamento progressivo. Due famiglie di soluzioni lo affrontano in modo diverso.

\begin{itemize}
  \item Le architetture \textbf{encoder-decoder} con connessioni di salto, come
  U-Net~\cite{ronneberger2015unet}, recuperano il dettaglio combinando, a ogni scala del
  decoder, le mappe ad alta risoluzione dell'encoder. Con un encoder residuale~\cite{he2016resnet}
  si ottengono le varianti \emph{ResU-Net}, che possono riusare pesi pre-addestrati su
  ImageNet~\cite{deng2009imagenet}.
  \item Le architetture con \textbf{convoluzioni dilatate} (\emph{atrous}), come
  DeepLab~\cite{chen2018deeplab}, allargano il campo recettivo senza ridurre la risoluzione,
  a un costo di memoria maggiore.
\end{itemize}

Più di recente gli encoder di tipo transformer, come ViT~\cite{dosovitskiy2021vit} e Swin
Transformer~\cite{liu2021swin}, e i decoder leggeri come quello di
SegFormer~\cite{xie2021segformer} hanno spostato lo stato dell'arte, al prezzo di un costo
computazionale e di un fabbisogno di dati maggiori.

\section{Segmentazione di immagini aeree e il benchmark ISPRS}
\label{sec:stato-aereo}

Il benchmark ISPRS~\cite{rottensteiner2012isprs} ha reso confrontabili i metodi di
classificazione di scene urbane aeree. Tra i primi lavori con reti profonde,
Sherrah~\cite{sherrah2016fcn} e Marmanis et al.~\cite{marmanis2016ensemble} hanno applicato FCN
e loro ensemble. Audebert et al.~\cite{audebert2018beyond} hanno mostrato il vantaggio di
fondere l'informazione altimetrica (DSM/nDSM) con quella spettrale, sia con una fusione precoce
(come canale aggiuntivo) sia con una fusione tardiva tra due rami. Lo stesso gruppo di autori di
\maresunet{} ha poi proposto UNetFormer~\cite{wang2022unetformer}, un decoder con attenzione
globale-locale per scene urbane. MMSegmentation~\cite{mmseg2020} ha consolidato uno split di
riferimento per Potsdam e Vaihingen, adottato da gran parte della letteratura recente, che questa
tesi segue (\cref{sec:split}).

Due caratteristiche distinguono il dominio dalle immagini naturali e guidano le scelte di
questo lavoro. La prima è l'\textbf{assenza di un'orientazione canonica}: una scena aerea
ruotata di \ang{90} o riflessa resta una scena plausibile, quindi le trasformazioni del gruppo
diedrale $D_4$ sono augmentation \enquote{gratuite}, che non alterano le etichette. La seconda è
la \textbf{disponibilità di modalità ausiliarie} co-registrate, come il DSM e l'nDSM, che
portano un'informazione (l'altezza) non ricavabile dal solo colore.

\section{Attenzione efficiente}
\label{sec:stato-attenzione}

L'attenzione \emph{dot-product}~\cite{vaswani2017attention} pesa ogni posizione in base alla
somiglianza tra la sua \emph{query} e le \emph{key} di tutte le altre. Nella visione artificiale è stata
introdotta come blocco \emph{non-local}~\cite{wang2018nonlocal} e, per la segmentazione, da
DANet~\cite{fu2019danet}, che affianca un'attenzione sulle posizioni (PAM) a una sui canali
(CAM). Il costo di memoria e di calcolo cresce però con il quadrato del numero di posizioni
$N=HW$, e questo ne limita l'uso alle mappe a bassa risoluzione.

Le \emph{attenzioni lineari} riscrivono il prodotto $\softmax(QK\T)V$ come $\phi(Q)\bigl(\phi(K)\T
V\bigr)$ per un'opportuna mappa di caratteristiche $\phi$, e sfruttano l'associatività del
prodotto matriciale per ridurre il costo a $\bigO{N}$~\cite{katharopoulos2020linear}. Il
\emph{Linear Attention Mechanism} (LAM) di Li et al.~\cite{li2020lam} ottiene lo stesso risultato
con uno sviluppo di Taylor al primo ordine dell'esponenziale, applicato a query e key
normalizzate. \maresunet{}~\cite{li2022maresunet} inserisce un blocco PAM lineare + CAM in ogni
connessione di salto di una ResU-Net. La derivazione completa è nella
\cref{sec:attenzione-lineare}.

\section{Rappresentazioni discrete}
\label{sec:stato-discrete}

Gli autoencoder variazionali~\cite{kingma2014vae} apprendono una rappresentazione latente
continua regolarizzata verso una distribuzione a priori. Il VQ-VAE~\cite{vandenoord2017vqvae}
sostituisce lo spazio latente continuo con un insieme finito di vettori (\emph{codebook}): ogni
vettore prodotto dall'encoder viene sostituito dal codice più vicino. L'operazione non è
derivabile e il gradiente viene propagato con lo stimatore
\emph{straight-through}~\cite{bengio2013straight}. Le rappresentazioni discrete sono state alla
base di modelli generativi ad alta fedeltà, come VQ-VAE-2~\cite{razavi2019vqvae2} e
VQGAN~\cite{esser2021taming}. Il loro problema pratico più noto è il \emph{collasso del
codebook}: pochi codici vengono usati e gli altri restano \enquote{morti}. Le contromisure più
diffuse sono l'aggiornamento per medie mobili esponenziali (EMA), l'inizializzazione dai dati e
il riavvio dei codici inutilizzati~\cite{lancucki2020robust}. Per la segmentazione, il codebook
agisce come un \emph{information bottleneck}~\cite{tishby2000ib}: limita la quantità di
informazione che attraversa il collo di bottiglia, con un possibile effetto regolarizzante.
L'idea richiama i \emph{visual words} di OBoW~\cite{gidaris2021obow}, in cui le feature vengono
assegnate alla parola più vicina di un vocabolario.

\section{Apprendimento auto-supervisionato}
\label{sec:stato-ssl}

L'apprendimento auto-supervisionato (\emph{self-supervised learning}, SSL) costruisce un compito
ausiliario a partire dai soli dati, senza etichette, e usa l'encoder così addestrato come punto
di partenza per il compito a valle. La \cref{tab:ssl-famiglie} riassume le famiglie principali
nel contesto di questa tesi. La sfida comune a tutti i metodi che confrontano due viste della
stessa immagine è il \emph{collasso}: la soluzione banale in cui la rete produce lo stesso
output per ogni ingresso.

\begin{table}[t]
  \centering
  \small
  \caption[Famiglie di metodi auto-supervisionati]{Famiglie di metodi di apprendimento
    auto-supervisionato per immagini e meccanismo con cui ciascuna evita il collasso.}
  \label{tab:ssl-famiglie}
  \begin{tabular}{|l|p{3.1cm}|p{3.7cm}|l|}
    \hline
    \textbf{Famiglia} & \textbf{Compito} & \textbf{Contro il collasso} & \textbf{Esempi} \\
    \hline
    Contrastiva & distinguere viste della stessa immagine da quelle di altre &
      esempi negativi (batch o coda) & \cite{chen2020simclr,he2020moco} \\
    \hline
    Clustering & predire l'assegnazione a prototipi di un'altra vista &
      vincolo di equipartizione & \cite{caron2020swav} \\
    \hline
    Non contrastiva & predire la rappresentazione di un'altra vista &
      asimmetria: predittore e \emph{stop-gradient}, teacher EMA & \cite{grill2020byol,chen2021simsiam} \\
    \hline
    Auto-distillazione & lo student imita la distribuzione del teacher &
      \emph{centering} e \emph{sharpening} dell'output del teacher, teacher EMA &
      \cite{caron2021dino,li2022esvit} \\
    \hline
    \emph{Bag of visual words} & ricostruire l'istogramma di parole visive del teacher &
      vocabolario aggiornato online, teacher EMA & \cite{gidaris2021obow} \\
    \hline
    Ricostruzione mascherata & ricostruire le porzioni mascherate &
      compito generativo, nessuna coppia di viste & \cite{he2022mae,cong2022satmae} \\
    \hline
  \end{tabular}
\end{table}

L'auto-distillazione discende dalla distillazione della conoscenza di Hinton et
al.~\cite{hinton2015distilling,gou2021kdsurvey}. Il teacher non è più una rete pre-addestrata
più grande, ma una media mobile esponenziale dei pesi dello student, come nel \emph{mean
teacher}~\cite{tarvainen2017meanteacher} e nella media di Polyak~\cite{polyak1992averaging}.
DINO~\cite{caron2021dino} combina questa idea con un'augmentation \emph{multi-crop} (due viste
globali per il teacher, viste globali e locali per lo student) e con la coppia
\emph{centering}/\emph{sharpening}. È stato proposto per i Vision Transformer, ma funziona anche
con encoder ResNet. Le versioni successive, come DINOv2~\cite{oquab2024dinov2}, ne hanno
confermato la scalabilità a grandi collezioni di immagini.

Nel telerilevamento l'SSL è motivato proprio dalla scarsità di etichette~\cite{wang2022ssl_rs}.
Metodi specifici sfruttano la dimensione temporale (SeCo~\cite{manas2021seco}), la
geolocalizzazione (GASSL~\cite{ayush2021gassl}) o la ricostruzione mascherata di immagini
multispettrali (SatMAE~\cite{cong2022satmae}). MARE~\cite{marsocci2021mare} ha applicato OBoW
alla \maresunet{} su Vaihingen e ha osservato che le augmentation cromatiche forti, utili su
ImageNet, non aiutano necessariamente sulle immagini aeree. È l'osservazione all'origine della
\RQ{4}.

\section{La tesi del 2022}
\label{sec:tesi-2022}

La tesi del 2022~\cite{campana2022tesi} ha studiato una griglia di $2\times2\times2$
configurazioni: architettura (\maresunet{}, \maresunetvq{}), encoder (ResNet-18, ResNet-34) e
pre-addestramento (ImageNet supervisionato, DINO su Potsdam). Le scelte sperimentali principali
erano le seguenti.

\begin{itemize}
  \item \textbf{Dati}: Potsdam, ritagli di $256\times256$ non sovrapposti. Erano state rimosse
  dieci aree giudicate troppo omogenee (\tile{2\_10}, \tile{2\_11}, \tile{2\_12}, \tile{2\_13},
  \tile{2\_14}, \tile{3\_10}, \tile{6\_7}, \tile{6\_13}, \tile{7\_13}, \tile{7\_9}); delle 28
  rimanenti, 22 erano usate per il training e 6 per la validation (circa \num{11000} ritagli di
  training).
  \item \textbf{Addestramento}: Adam, learning rate $10^{-3}$, StepLR (passo 15, $\gamma=0{,}4$),
  100 epoche, batch 32, nessuna augmentation in fine-tuning.
  \item \textbf{Pre-addestramento}: DINO con scheduler coseno, \emph{weight decay} da 0,04 a 0,4,
  $\tau_s=0{,}1$ e $\tau_t$ da 0,04 a 0,07.
\end{itemize}

\begin{table}[t]
  \centering
  \small
  \caption[Risultati di test della tesi del 2022]{Risultati di test riportati nella tesi del
    2022 su Potsdam (tab.~4.5 dell'originale), con il protocollo di allora. In grassetto il
    valore migliore di ogni colonna.}
  \label{tab:risultati-2022}
  \begin{tabular}{|l|l|l|c|c|c|}
    \hline
    \textbf{Modello} & \textbf{Encoder} & \textbf{Pre-addestr.} & \textbf{OA} & \textbf{mF1} & \textbf{mIoU} \\
    \hline
    \maresunet{}   & ResNet-18 & ImageNet & 0,791 & \textbf{0,722} & \textbf{0,604} \\
    \maresunet{}   & ResNet-34 & ImageNet & \textbf{0,791} & 0,722 & 0,604 \\
    \maresunet{}   & ResNet-18 & DINO     & 0,741 & 0,657 & 0,530 \\
    \maresunet{}   & ResNet-34 & DINO     & 0,739 & 0,653 & 0,525 \\
    \maresunetvq{} & ResNet-18 & ImageNet & 0,750 & 0,668 & 0,543 \\
    \maresunetvq{} & ResNet-34 & ImageNet & 0,744 & 0,660 & 0,534 \\
    \maresunetvq{} & ResNet-18 & DINO     & 0,702 & 0,593 & 0,465 \\
    \maresunetvq{} & ResNet-34 & DINO     & 0,703 & 0,595 & 0,467 \\
    \hline
  \end{tabular}
\end{table}

I risultati (\cref{tab:risultati-2022}) indicavano che la quantizzazione e il pre-addestramento
DINO peggioravano tutte le metriche e che i loro effetti si sommavano. ResNet-18 e ResNet-34
risultavano equivalenti. Il clutter era la classe peggiore (IoU $\approx 0{,}15$) e alberi e
vegetazione bassa si confondevano. Le metriche di validation (mIoU $\approx 0{,}96$) erano
molto più alte di quelle di test ($\approx 0{,}60$). La tesi attribuiva il divario
all'overfitting e la sconfitta di DINO alla scarsità di dati di pre-addestramento (circa
\num{11000} ritagli contro i milioni di immagini di ImageNet).

\section{Lettura critica}
\label{sec:lettura-critica}

Il confronto tra il testo, il codice originale e i numeri riportati fa emergere diversi
fattori, \emph{alternativi} alle spiegazioni proposte, che possono aver prodotto il risultato
negativo (\cref{tab:lettura-critica}). Il punto per questo lavoro è che le ipotesi di ricerca
\textbf{non sono state confutate in modo conclusivo}.

\begin{table}[t]
  \centering
  \small
  \caption[Osservazioni della tesi del 2022 e fattori alternativi]{Osservazioni della tesi del
    2022, spiegazioni proposte allora e fattori alternativi individuati nel codice e nel
    protocollo.}
  \label{tab:lettura-critica}
  \begin{tabular}{|p{2.6cm}|p{2.6cm}|p{7.2cm}|}
    \hline
    \textbf{Osservazione} & \textbf{Spiegazione proposta} & \textbf{Fattore alternativo} \\
    \hline
    Clutter quasi mai riconosciuto & eterogeneità della classe &
      La perdita escludeva l'indice 0, che corrispondeva proprio al clutter: per quella classe
      il modello riceveva solo gradienti che ne abbassavano il logit
      (\cref{prop:classe-esclusa}). \\
    \hline
    Validation $\gg$ test & overfitting &
      L'overfitting si manifesta come training $\gg$ validation, non validation $\gg$ test. Un
      salto da 0,96 a 0,60 suggerisce una validation non indipendente (ritagli sovrapposti o
      split non disgiunti per area, \cite{kaufman2012leakage}) o metriche calcolate in modo
      diverso (media per batch, \cref{sec:metriche}). \\
    \hline
    DINO $<$ ImageNet & dati di pre-addestramento insufficienti &
      L'implementazione si discostava da DINO: nessuna \emph{projection head}, softmax
      calcolata sull'intero batch appiattito (le distribuzioni dei campioni si mescolano),
      viste locali a 224 invece che a 96 pixel, temperatura del teacher fissa, nessuno schedule
      del \emph{weight decay}, teacher non inizializzato come copia dello student. \\
    \hline
    VQ peggiora & la quantizzazione non aiuta &
      Il codebook quantizzava solo il collo di bottiglia, aggirato dalle connessioni di salto
      (\cref{prop:aggiramento}). Lo stato EMA veniva ricreato a ogni passo. Un encoder
      condiviso per errore tra istanze del modello faceva sì che congelare il teacher
      congelasse anche lo student. \\
    \hline
    ResNet-18 $\approx$ ResNet-34 & basta il modello piccolo &
      Nel codice disponibile non c'è una variante ResNet-34: il codice che ha prodotto le
      tabelle non è con certezza quello conservato. \\
    \hline
  \end{tabular}
\end{table}

A questi fattori si aggiungono alcune scelte di protocollo che impediscono il confronto con la
letteratura.

\begin{itemize}
  \item \textbf{Split non standard}: quattro delle aree rimosse (\tile{2\_13}, \tile{2\_14},
  \tile{6\_13}, \tile{7\_13}) sono tile di \emph{test} nello split più usato, e il testo non
  descrive da dove provengano i \enquote{Test Scores}.
  \item \textbf{Ingressi non normalizzati} con le statistiche di ImageNet, pur partendo da
  pesi ImageNet.
  \item \textbf{Batch di training ordinati} (nessun rimescolamento), quindi spazialmente
  correlati.
  \item \textbf{Una sola esecuzione per configurazione}, senza seed: differenze di uno o due
  punti non sono distinguibili dalla variabilità tra esecuzioni~\cite{bouthillier2021variance}.
\end{itemize}

\section{Implicazioni per questo lavoro}
\label{sec:implicazioni}

La \cref{tab:correzioni} riassume come ciascun problema viene affrontato nella nuova
implementazione e dove la scelta è discussa. Le correzioni non cambiano le domande di ricerca:
le rendono rispondibili.

\begin{table}[t]
  \centering
  \small
  \caption[Problemi del protocollo 2022 e correzioni adottate]{Problemi del protocollo 2022 e
    correzioni adottate in questo lavoro.}
  \label{tab:correzioni}
  \begin{tabular}{|p{4.4cm}|p{6.0cm}|c|}
    \hline
    \textbf{Problema} & \textbf{Correzione} & \textbf{Dove} \\
    \hline
    Clutter escluso dalla perdita & si ignora solo il valore 255 (pixel senza etichetta) & \cref{sec:formulazione} \\
    \hline
    Metriche mediate per batch & matrice di confusione globale sull'intero split & \cref{sec:metriche} \\
    \hline
    Validation non indipendente & split per tile intere; validation su tile escluse dal training & \cref{sec:split} \\
    \hline
    Split non standard & split di riferimento della letteratura; protocollo 2022 mantenuto solo per il confronto & \cref{sec:split} \\
    \hline
    Valutazione su ritagli & finestra scorrevole su tile intere, etichette piene ed \emph{eroded} & \cref{sec:finestra} \\
    \hline
    DINO non fedele & head di proiezione, perdita per campione, multi-crop $2\times224 + 6\times96$, schedule completi & \cref{sec:dino} \\
    \hline
    Codebook fragile & EMA in variabili persistenti, inizializzazione dai dati, riavvio dei codici morti, monitoraggio della perplexity & \cref{sec:vq} \\
    \hline
    Collo di bottiglia aggirato & quantizzazione opzionale anche delle connessioni di salto & \cref{sec:vq-aggiramento} \\
    \hline
    Normalizzazione assente & normalizzazione ImageNet come primo strato del modello & \cref{sec:addestramento} \\
    \hline
    Un'esecuzione, nessun seed & tre seed per configurazione, media $\pm$ deviazione standard & \cref{sec:statistica} \\
    \hline
  \end{tabular}
\end{table}

% =====================================================================
%  Capitolo 3 — Fondamenti teorici
% =====================================================================
\chapter{Fondamenti teorici}
\label{cap:teoria}

Questo capitolo sviluppa il quadro teorico dei meccanismi studiati. Si parte dalla formulazione
del problema e dalla funzione di perdita, poi si passa all'architettura encoder-decoder e ai
meccanismi di attenzione, fino alla derivazione dell'attenzione lineare usata da \maresunet{}.
Seguono la quantizzazione vettoriale e l'auto-distillazione, la fusione dell'informazione
altimetrica, le metriche e l'inferenza su tile intere. Per ogni meccanismo si mettono in
evidenza le proprietà che motivano le domande di ricerca e le scelte del
\cref{cap:metodologia}.

% ---------------------------------------------------------------------
\section{Formulazione del problema}
\label{sec:formulazione}

Un'immagine è un tensore $\vect{x}\in\R^{H\times W\times C_{\mathrm{in}}}$, con $C_{\mathrm{in}}=3$
(RGB o IR-R-G) oppure $C_{\mathrm{in}}=4$ se si aggiunge l'nDSM. L'insieme delle posizioni è
$\Omega=\{1,\dots,H\}\times\{1,\dots,W\}$. L'etichetta è una mappa
$y:\Omega\to\{1,\dots,K\}\cup\{\varnothing\}$, dove $K=6$ è il numero di classi e
$\varnothing$ (codificato con il valore 255) indica un pixel senza etichetta valida. Le
posizioni etichettate sono $\Omega_v=\{u\in\Omega : y_u\neq\varnothing\}$.

Un modello di segmentazione $f_\theta$ associa a ogni posizione un vettore di logit
$\vect{z}_u=f_\theta(\vect{x})_u\in\R^K$ e quindi la distribuzione
$p_\theta(k\mid\vect{x},u)=\softmax(\vect{z}_u)_k$. L'addestramento minimizza l'entropia
incrociata media sulle posizioni etichettate:
\begin{equation}
  \mathcal{L}_{\mathrm{CE}}(\theta)
  = -\frac{1}{\abs{\Omega_v}}\sum_{u\in\Omega_v}\log p_\theta(y_u\mid\vect{x},u).
  \label{eq:ce-mascherata}
\end{equation}
La predizione è $\hat y_u=\argmax_k \vect{z}_{u,k}$. La scelta di \emph{quali} pixel entrano in
$\Omega_v$ non è un dettaglio implementativo, come mostra il risultato seguente.

\begin{proposizione}[Classe esclusa dalla perdita]
\label{prop:classe-esclusa}
Si escluda dalla perdita~\eqref{eq:ce-mascherata} una classe reale $c$, cioè si sostituisca
$\Omega_v$ con $\Omega_v'=\{u\in\Omega_v : y_u\neq c\}$. Allora, per ogni posizione
$u\in\Omega_v'$,
\[
  \frac{\partial\mathcal{L}}{\partial z_{u,c}}=\frac{p_\theta(c\mid\vect{x},u)}{\abs{\Omega_v'}}\;\geq\;0,
\]
mentre i pixel di classe $c$ non contribuiscono alla perdita. Ogni passo di discesa del gradiente
riduce quindi il logit della classe $c$ e non lo aumenta mai.
\end{proposizione}
\begin{proof}
Per una singola posizione $\ell_u=-\log\softmax(\vect{z}_u)_{y_u}$ si ha
$\partial\ell_u/\partial z_{u,k}=p_\theta(k\mid\vect{x},u)-\Ind[k=y_u]$. Per $u\in\Omega_v'$ vale
$y_u\neq c$, quindi la derivata rispetto a $z_{u,c}$ è $p_\theta(c\mid\vect{x},u)\ge0$. Le
posizioni con $y_u=c$ non compaiono nella somma.
\end{proof}

Nel codice del 2022 l'indice ignorato era 0, che corrispondeva proprio al clutter. Per la
\cref{prop:classe-esclusa} il modello riceveva solo una pressione a \emph{non} prevedere il
clutter. Poteva ancora prevederlo dove l'ingresso era diverso da tutti i pixel etichettati,
perché lì la pressione era solo indiretta: questo è compatibile con l'IoU basso ma non nullo
riportato allora ($\approx0{,}15$). In questo lavoro si ignora soltanto il valore $\varnothing$.

% ---------------------------------------------------------------------
\section{Encoder-decoder residuali}
\label{sec:encoder-decoder}

\paragraph{Blocchi residui.}
Un blocco residuo~\cite{he2016resnet} calcola $\vect{x}_{l+1}=\vect{x}_l+\mathcal{F}(\vect{x}_l;W_l)$.
Nella forma semplificata con mappa identità, srotolando la ricorrenza si ottiene
$\vect{x}_L=\vect{x}_l+\sum_{i=l}^{L-1}\mathcal{F}(\vect{x}_i;W_i)$ e quindi
\begin{equation}
  \frac{\partial\mathcal{L}}{\partial\vect{x}_l}
  =\frac{\partial\mathcal{L}}{\partial\vect{x}_L}
   \Bigl(I+\frac{\partial}{\partial\vect{x}_l}\sum_{i=l}^{L-1}\mathcal{F}(\vect{x}_i;W_i)\Bigr).
\end{equation}
Il termine identità fa arrivare il gradiente agli strati iniziali senza attenuazione
moltiplicativa. Per questo encoder profondi pre-addestrati si possono rifinire in modo stabile.

\paragraph{Gerarchia di scale.}
Gli encoder ResNet producono quattro mappe $\vect{E}_1,\dots,\vect{E}_4$ a stride $4,8,16,32$
(\cref{tab:stadi-resnet}). Le mappe superficiali sono ad alta risoluzione e semanticamente
povere, quelle profonde sono povere di dettaglio spaziale. Il numero di posizioni
$N_l=(H/s_l)(W/s_l)$ diminuisce di un fattore 4 a ogni stadio. È questa quantità a governare il
costo dell'attenzione (\cref{sec:attenzione-dot}).

\begin{table}[t]
  \centering
  \small
  \caption[Stadi degli encoder ResNet per un ritaglio di 512 pixel]{Stadi degli encoder ResNet
    per un ritaglio di $512\times512$ pixel: stride, canali, numero di posizioni $N_l$ e memoria
    della sola matrice di attenzione $N_l\times N_l$ in float32, per immagine e per testa.}
  \label{tab:stadi-resnet}
  \begin{tabular}{|c|c|c|c|c|r|}
    \hline
    \textbf{Stadio} & \textbf{Stride} & \makecell{\textbf{Canali}\\\textbf{R18/R34}} & \makecell{\textbf{Canali}\\\textbf{R50}} &
    \textbf{Posizioni $N_l$} & \makecell{\textbf{Memoria}\\\textbf{$N_l^2$}} \\
    \hline
    $\vect{E}_1$ & 4  & 64  & 256  & $128^2=\num{16384}$ & \SI{1}{GiB} \\
    $\vect{E}_2$ & 8  & 128 & 512  & $64^2=\num{4096}$   & \SI{64}{MiB} \\
    $\vect{E}_3$ & 16 & 256 & 1024 & $32^2=\num{1024}$   & \SI{4}{MiB} \\
    $\vect{E}_4$ & 32 & 512 & 2048 & $16^2=256$          & \SI{256}{KiB} \\
    \hline
  \end{tabular}
\end{table}

\paragraph{Decoder con connessioni di salto.}
Uno schema U-Net~\cite{ronneberger2015unet} ricostruisce la risoluzione combinando, a ogni scala,
l'uscita sovracampionata del livello inferiore con una trasformazione $\mathcal{S}_l$ della
mappa dell'encoder alla stessa scala:
\begin{equation}
  \vect{D}_l=\mathrm{Dec}_{l+1}(\vect{D}_{l+1})+\mathcal{S}_l(\vect{E}_l).
  \label{eq:decoder}
\end{equation}
Nella U-Net \enquote{semplice} $\mathcal{S}_l$ è l'identità. In \maresunet{} è un blocco di
attenzione (\cref{sec:maresunet}). È esattamente questa differenza a essere misurata dalla
\RQ{1}.

% ---------------------------------------------------------------------
\section{Meccanismi di attenzione}
\label{sec:attenzione}

\subsection{Attenzione dot-product}
\label{sec:attenzione-dot}

Sia $\vect{X}\in\R^{N\times C}$ una mappa di caratteristiche con le $N=hw$ posizioni disposte in
righe. L'attenzione dot-product~\cite{vaswani2017attention} proietta $\vect{X}$ in query, key e
value, $Q=\vect{X}W_q\in\R^{N\times d_k}$, $K=\vect{X}W_k\in\R^{N\times d_k}$ e
$V=\vect{X}W_v\in\R^{N\times C}$, e calcola
\begin{equation}
  D(Q,K,V)=\softmax_{\mathrm{riga}}\!\left(\frac{QK\T}{\sqrt{d_k}}\right)V .
  \label{eq:dot-product}
\end{equation}
L'uscita in posizione $i$ è una media pesata dei value di \emph{tutte} le posizioni, con pesi
dati dalla somiglianza tra la query $\vect{q}_i$ e le key $\vect{k}_j$. La matrice $QK\T$ ha
$N^2$ elementi: calcolo e memoria sono $\bigO{N^2 d_k}$ e $\bigO{N^2}$. Sul primo stadio di un
ritaglio di $512\times512$ pixel ($N=\num{16384}$) la sola matrice occupa \SI{1}{GiB} in
float32 per immagine e per testa (\cref{tab:stadi-resnet}). Con un batch di 8 immagini si
superano gli \SI{8}{\giga\byte} della GPU usata in questo lavoro prima ancora di contare le
attivazioni del resto della rete (\cref{fig:complessita}).

\begin{figure}[t]
  \centering
  % Memoria richiesta dall'attenzione sul primo livello dell'encoder (stride 4, C = 64,
% d_k = C/8) in funzione del lato del ritaglio, per una singola immagine in float32.
\begin{tikzpicture}
  \begin{axis}[
      width=0.9\linewidth, height=6.2cm,
      xlabel={lato del ritaglio $s$ (pixel)},
      ylabel={memoria (GiB, scala log.)},
      xmin=128, xmax=1024, ymode=log, ymin=1e-5, ymax=60,
      xtick={128,256,384,512,640,768,896,1024},
      grid=major, grid style={SapGrigioChiaro},
      legend pos=north west, legend cell align=left,
      legend style={font=\footnotesize, draw=SapGrigioTenue},
      tick label style={font=\footnotesize}, label style={font=\small},
      samples=60, domain=128:1024,
    ]
    % N = (s/4)^2 ; dot-product: N^2 * 4 B ; LAM: (N (C + d_k) + d_k C) * 4 B
    \addplot[SapRosso, very thick] {((x/4)^2)^2*4/2^30};
    \addlegendentry{dot-product: matrice $N\times N$}
    \addplot[SapGrigio, very thick, dashed] {(((x/4)^2)*(64+8) + 8*64)*4/2^30};
    \addlegendentry{LAM: stato $d_k\times C$ e uscite $N\times(C+d_k)$}
    \addplot[black, thin, densely dotted, domain=128:1024] {8};
    \node[font=\scriptsize, anchor=south east] at (axis cs:1020,8) {\SI{8}{\giga\byte} (RTX 4060)};
    \addplot[only marks, mark=*, mark size=1.6pt, SapRossoScuro] coordinates {(512,1)};
    \node[font=\scriptsize, anchor=west, text=SapRossoScuro] at (axis cs:525,0.6) {\SI{1}{GiB} per immagine a $s=512$};
  \end{axis}
\end{tikzpicture}

  \caption[Memoria dell'attenzione dot-product e lineare]{Memoria richiesta dall'attenzione sul
    primo stadio dell'encoder (stride 4, $C=64$, $d_k=C/8$) al variare del lato del ritaglio,
    per una singola immagine in float32. La matrice $N\times N$ cresce con la quarta potenza
    del lato. Lo stato dell'attenzione lineare cresce con il quadrato, cioè linearmente in
    $N$.}
  \label{fig:complessita}
\end{figure}

\subsection{Attenzione come media a nucleo}
\label{sec:attenzione-nucleo}

La riga $i$ della~\eqref{eq:dot-product} si può scrivere come
\begin{equation}
  D_i=\frac{\sum_{j=1}^{N}\mathrm{sim}(\vect{q}_i,\vect{k}_j)\,\vect{v}_j}
           {\sum_{j=1}^{N}\mathrm{sim}(\vect{q}_i,\vect{k}_j)},
  \qquad \mathrm{sim}(\vect{q},\vect{k})=\exp\!\bigl(\vect{q}\T\vect{k}/\sqrt{d_k}\bigr),
  \label{eq:attenzione-nucleo}
\end{equation}
cioè come una media a nucleo dei value. La~\eqref{eq:attenzione-nucleo} definisce
un'attenzione valida per qualunque funzione di somiglianza non negativa. Se
$\mathrm{sim}(\vect{q},\vect{k})=\phi(\vect{q})\T\psi(\vect{k})$ per due mappe
$\phi,\psi:\R^{d_k}\to\R^{m}$, la proprietà associativa del prodotto
matriciale~\cite{katharopoulos2020linear} consente di raccogliere le somme sulle key, che non
dipendono da $i$:
\begin{equation}
  D_i=\frac{\phi(\vect{q}_i)\T\Bigl(\sum_{j}\psi(\vect{k}_j)\,\vect{v}_j\T\Bigr)}
           {\phi(\vect{q}_i)\T\sum_{j}\psi(\vect{k}_j)} .
  \label{eq:attenzione-lineare-generale}
\end{equation}
La matrice $\sum_j\psi(\vect{k}_j)\vect{v}_j\T\in\R^{m\times C}$ e il vettore
$\sum_j\psi(\vect{k}_j)\in\R^{m}$ si calcolano una sola volta in $\bigO{NmC}$. Ogni uscita
costa poi $\bigO{mC}$: il costo complessivo è lineare in $N$ (\cref{fig:associativita}).

\begin{figure}[t]
  \centering
  \resizebox{0.92\linewidth}{!}{% Ordine dei prodotti: (Q K^T) V costa O(N^2 d), Q (K^T V) costa O(N d C).
\begin{tikzpicture}[x=1mm, y=1mm, font=\footnotesize]
  \tikzset{mat/.style={draw=SapRossoScuro, line width=0.5pt}}
  % --- sinistra: (Q K^T) V -------------------------------------------------
  \begin{scope}
    \draw[mat, fill=SapRossoChiaro] (0,0) rectangle (4,24);      \node at (2,-3) {$Q$};
    \node at (6,12) {$\cdot$};
    \draw[mat, fill=SapRossoChiaro] (8,20) rectangle (32,24);    \node at (20,17) {$K\T$};
    \node at (36,12) {$=$};
    \draw[mat, fill=SapRossoTenue] (40,0) rectangle (64,24);     \node at (52,12) {$N\times N$};
    \node at (68,12) {$\cdot$};
    \draw[mat, fill=SapGrigioChiaro, draw=SapGrigio] (72,0) rectangle (78,24); \node at (75,-3) {$V$};
    \node[text=SapRossoScuro, font=\footnotesize\bfseries] at (39,31) {$(QK\T)\,V$: memoria e calcolo $\mathcal{O}(N^2)$};
  \end{scope}
  % --- destra: Q (K^T V) -----------------------------------------------------
  \begin{scope}[xshift=92mm]
    \draw[mat, fill=SapRossoChiaro] (0,20) rectangle (24,24);     \node at (12,17) {$\hat K\T$};
    \node at (28,12) {$\cdot$};
    \draw[mat, fill=SapGrigioChiaro, draw=SapGrigio] (32,0) rectangle (38,24); \node at (35,-3) {$V$};
    \node at (42,12) {$=$};
    \draw[mat, fill=SapRossoTenue] (46,18) rectangle (52,24);   \node[align=center] at (49,13) {$d_k\times C$};
    \node at (57,12) {$\Rightarrow$};
    \draw[mat, fill=SapRossoChiaro] (62,0) rectangle (66,24);   \node at (64,-3) {$\hat Q$};
    \node at (69,12) {$\cdot$};
    \draw[mat, fill=SapRossoTenue] (72,18) rectangle (78,24);
    \node[text=SapRossoScuro, font=\footnotesize\bfseries] at (39,31) {$\hat Q\,(\hat K\T V)$: memoria e calcolo $\mathcal{O}(N)$};
  \end{scope}
\end{tikzpicture}
}
  \caption[Ordine dei prodotti nell'attenzione]{L'ordine dei prodotti determina il costo. A
    sinistra la forma dot-product materializza la matrice $N\times N$. A destra la forma lineare
    calcola prima $\hat K\T V\in\R^{d_k\times C}$, comune a tutte le query.}
  \label{fig:associativita}
\end{figure}

\subsection{Il Linear Attention Mechanism}
\label{sec:attenzione-lineare}

Il LAM di Li et al.~\cite{li2020lam} ottiene una fattorizzazione del tipo
\eqref{eq:attenzione-lineare-generale} sviluppando l'esponenziale al primo ordine,
$e^{\vect{q}\T\vect{k}}\approx1+\vect{q}\T\vect{k}$. Lo sviluppo da solo non garantisce pesi non
negativi, perché $\vect{q}\T\vect{k}$ può essere minore di $-1$. Per questo query e key vengono
normalizzate in norma $\ell_2$, $\hatq_i=\vect{q}_i/\norm{\vect{q}_i}_2$ e
$\hatk_j=\vect{k}_j/\norm{\vect{k}_j}_2$, e si usa
\begin{equation}
  \mathrm{sim}(\vect{q}_i,\vect{k}_j)=1+\hatq_i\T\hatk_j=1+\cos\vartheta_{ij}\in[0,2].
\end{equation}
Sostituendo nella~\eqref{eq:attenzione-nucleo} si ottiene
\begin{equation}
  D_i=\frac{\sum_{j}\vect{v}_j+\Bigl(\sum_{j}\hatk_j\vect{v}_j\T\Bigr)\T\hatq_i}
           {N+\hatq_i\T\sum_{j}\hatk_j},
  \label{eq:lam}
\end{equation}
che corrisponde alla~\eqref{eq:attenzione-lineare-generale} con
$\phi(\vect{q})=(1,\hat{\vect{q}})$ e $\psi(\vect{k})=(1,\hat{\vect{k}})$, quindi con $m=d_k+1$.
In forma matriciale, con $\hat Q$ e $\hat K$ normalizzate riga per riga,
$D=\bigl(\vect{1}\vect{1}\T V+\hat Q(\hat K\T V)\bigr)\oslash\bigl(N\vect{1}+\hat Q\hat K\T\vect{1}\bigr)$,
dove $\oslash$ divide ogni riga per il corrispondente elemento del vettore. Il calcolo richiede
$\bigO{N d_k C}$ operazioni e memoria $\bigO{N(d_k+C)+d_kC}$. Nell'implementazione si aggiunge
una piccola costante $\epsilon$ alla somma delle key per stabilità numerica.

\begin{osservazione}[Contrasto limitato del LAM]
\label{oss:contrasto}
Sia $\bar{\vect{m}}=\frac1N\sum_j\hatk_j$ la media delle key normalizzate, con
$\norm{\bar{\vect{m}}}_2\le1$. Il peso che la posizione $i$ assegna alla posizione $j$ è
\[
  \alpha_{ij}=\frac{1+\hatq_i\T\hatk_j}{N\,(1+\hatq_i\T\bar{\vect{m}})}
  \;\leq\;\frac{2}{N\,(1-\norm{\bar{\vect{m}}}_2)} .
\]
Se le key non sono tutte allineate ($\norm{\bar{\vect{m}}}_2$ lontana da 1), nessuna posizione
riceve più di una piccola frazione costante del peso medio $1/N$. La softmax, al contrario, può
concentrare quasi tutto il peso su una sola posizione. Il LAM si comporta quindi come
un'\emph{aggregazione di contesto globale modulata} dalla somiglianza coseno più che come un
meccanismo di selezione. Il caso $\norm{\bar{\vect{m}}}_2\to1$ con
$\hatq_i\approx-\bar{\vect{m}}$ è quello in cui il denominatore si annulla e serve
$\epsilon$.
\end{osservazione}

L'osservazione suggerisce che il contributo del LAM vada cercato nella coerenza su larga scala
(per esempio nell'uniformità delle grandi superfici) più che nella precisione dei contorni.
\Segnaposto{Verificare l'osservazione sui modelli addestrati: distribuzione di
$\norm{\bar{\vect{m}}}_2$ e di $\max_j\alpha_{ij}$ per stadio, da confrontare con le classi in
cui \RQ{1} mostra differenze.}

L'uscita del blocco è residuale, $\vect{X}'=\vect{X}+\gamma\,D(\vect{X})$, con un guadagno
scalare $\gamma$ appreso e \textbf{inizializzato a zero}. All'inizio dell'addestramento il
blocco è quindi l'identità e il modello parte dal comportamento della rete senza attenzione.

\subsection{Attenzione sui canali}
\label{sec:cam}

Sui canali il costo quadratico non è un problema, perché $C\ll N$. Seguendo
DANet~\cite{fu2019danet}, il modulo CAM calcola la matrice di affinità tra canali
$E=\vect{X}\T\vect{X}\in\R^{C\times C}$ e i pesi
$A=\softmax_{\mathrm{riga}}\bigl(\max_{\mathrm{riga}}(E)-E\bigr)$, poi produce
$\vect{X}'=\vect{X}+\beta\,\vect{X}A\T$ con $\beta$ inizializzato a zero. Il costo è
$\bigO{NC^2}$.

\begin{osservazione}
La softmax è invariante per l'aggiunta di una costante a ciascuna riga, quindi
$\softmax_{\mathrm{riga}}(\max_{\mathrm{riga}}(E)-E)=\softmax_{\mathrm{riga}}(-E)$. La
sottrazione del massimo è solo un accorgimento di stabilità numerica, e l'effetto netto è di
dare più peso ai canali \emph{meno} correlati con il canale considerato.
\end{osservazione}

\subsection{L'architettura MAResU-Net}
\label{sec:maresunet}

\maresunet{}~\cite{li2022maresunet} instanzia la~\eqref{eq:decoder} con un blocco di attenzione
su ciascuna delle quattro scale (\cref{fig:maresunet,fig:pamcam}). Il blocco applica una
convoluzione $3\times3$, poi in parallelo il PAM lineare~\eqref{eq:lam} e il CAM, ciascuno
seguito da dropout spaziale e convoluzione $1\times1$. Le due uscite vengono sommate e
proiettate con una convoluzione $1\times1$. Il decoder è
\begin{align}
  \vect{S}_l &= \mathrm{PAMCAM}_l(\vect{E}_l),\quad l=1,\dots,4, &
  \vect{B} &= \mathrm{Neck}(\vect{S}_4), \nonumber\\
  \vect{D}_3 &= \mathrm{Dec}_4(\vect{B})+\vect{S}_3, &
  \vect{D}_2 &= \mathrm{Dec}_3(\vect{D}_3)+\vect{S}_2, \label{eq:maresunet}\\
  \vect{D}_1 &= \mathrm{Dec}_2(\vect{D}_2)+\vect{S}_1, &
  \vect{z} &= \mathrm{Head}\bigl(\mathrm{Dec}_1(\vect{D}_1)\bigr). \nonumber
\end{align}
Ogni $\mathrm{Dec}_l$ è la sequenza convoluzione $1\times1$ (riduzione dei canali a un quarto),
convoluzione trasposta $3\times3$ con stride 2 e convoluzione $1\times1$, ciascuna con batch
normalization~\cite{ioffe2015batchnorm} e ReLU. La testa riporta la risoluzione piena con una
convoluzione trasposta $4\times4$ con stride 2 seguita da due convoluzioni $3\times3$. Il
\emph{neck} è l'identità nella \maresunet{} e un quantizzatore vettoriale nella \maresunetvq{}
(\cref{sec:vq}).

\begin{figure}[p]
  \centering
  \resizebox{\linewidth}{!}{% Architettura MAResU-Net: encoder ResNet, blocchi PAM-CAM sulle connessioni di salto,
% neck opzionale (identità o VQ), decoder con somme e testa a risoluzione piena.
\begin{tikzpicture}[
    e/.style={enc, minimum width=2.3cm},
    a/.style={attn, minimum width=2.3cm},
    d/.style={dec, minimum width=2.3cm},
    somma/.style={circle, draw=SapRossoScuro, fill=white, inner sep=0pt, minimum size=4.2mm,
                  font=\footnotesize, line width=0.5pt},
  ]
  % --- livelli (stride) ---------------------------------------------------
  \def\yA{-2.0}\def\yB{-3.6}\def\yC{-5.2}\def\yD{-6.8}

  % --- encoder --------------------------------------------------------------
  \node[dato, minimum width=2.3cm] (in) at (0,0.1) {Immagine\\{\scriptsize $H\times W\times3$ (o 4 con nDSM)}};
  \node[e] (stem) at (0,-0.95) {conv $7{\times}7$, s2 + max-pool};
  \node[e] (e1) at (0,\yA) {\texttt{layer1}\\{\scriptsize 64 canali}};
  \node[e] (e2) at (0,\yB) {\texttt{layer2}\\{\scriptsize 128 canali}};
  \node[e] (e3) at (0,\yC) {\texttt{layer3}\\{\scriptsize 256 canali}};
  \node[e] (e4) at (0,\yD) {\texttt{layer4}\\{\scriptsize 512 canali}};
  \draw[flusso] (in) -- (stem);
  \draw[flusso] (stem) -- (e1);
  \draw[flusso] (e1) -- (e2);
  \draw[flusso] (e2) -- (e3);
  \draw[flusso] (e3) -- (e4);

  % --- etichette di stride ----------------------------------------------------
  \foreach \y/\s in {\yA/{1/4}, \yB/{1/8}, \yC/{1/16}, \yD/{1/32}}{
    \node[nota, anchor=east] at (-1.45,\y) {$\s$};
  }
  \node[nota, anchor=east, font=\scriptsize\bfseries] at (-1.45,-1.45) {scala};

  % --- attenzione sulle connessioni di salto ---------------------------------
  \node[a] (a1) at (3.9,\yA) {PAM lineare $\parallel$ CAM};
  \node[a] (a2) at (3.9,\yB) {PAM lineare $\parallel$ CAM};
  \node[a] (a3) at (3.9,\yC) {PAM lineare $\parallel$ CAM};
  \node[a] (a4) at (3.9,\yD) {PAM lineare $\parallel$ CAM};
  \foreach \i in {1,...,4}{ \draw[flusso] (e\i) -- (a\i); }

  % --- neck ------------------------------------------------------------------
  \node[quant, minimum width=1.9cm] (neck) at (6.9,\yD) {Neck\\{\scriptsize identità o VQ}};
  \draw[flusso] (a4) -- (neck);

  % --- decoder ----------------------------------------------------------------
  \node[d] (d4) at (9.6,{(\yC+\yD)/2}) {Dec$_4$ \scriptsize($\uparrow2$)};
  \node[somma] (s3) at (9.6,\yC) {$+$};
  \node[d] (d3) at (9.6,{(\yB+\yC)/2}) {Dec$_3$ \scriptsize($\uparrow2$)};
  \node[somma] (s2) at (9.6,\yB) {$+$};
  \node[d] (d2) at (9.6,{(\yA+\yB)/2}) {Dec$_2$ \scriptsize($\uparrow2$)};
  \node[somma] (s1) at (9.6,\yA) {$+$};
  \node[d] (d1) at (9.6,-0.95) {Dec$_1$ \scriptsize($\uparrow2$, 1/2)};
  \node[d, fill=SapRossoVelo, text width=2.5cm] (head) at (9.6,0.25) {Testa\\{\scriptsize ConvT $4{\times}4$ s2,\\2 conv $3{\times}3$}};
  \node[dato, anchor=west, minimum width=1.6cm] (out) at (11.8,0.25) {Logit\\{\scriptsize $H\times W\times6$}};

  \draw[flusso] (neck) -| (d4);
  \draw[flusso] (d4) -- (s3);
  \draw[flusso] (s3) -- (d3);
  \draw[flusso] (d3) -- (s2);
  \draw[flusso] (s2) -- (d2);
  \draw[flusso] (d2) -- (s1);
  \draw[flusso] (s1) -- (d1);
  \draw[flusso] (d1) -- (head);
  \draw[flusso] (head) -- (out);

  % connessioni di salto (dopo l'attenzione)
  \draw[skipc] (a3) -- (s3);
  \draw[skipc] (a2) -- (s2);
  \draw[skipc] (a1) -- (s1);

  % --- raggruppamenti -----------------------------------------------------------
  \begin{scope}[on background layer]
    \node[gruppo, fit=(stem)(e4), label={[nota]below:encoder (ResNet-18/34)}] {};
    \node[gruppo, fit=(d1)(d4)(head), label={[nota]above:decoder}] {};
  \end{scope}
\end{tikzpicture}
}
  \caption[Architettura MAResU-Net]{Architettura \maresunet{}. Le mappe dell'encoder a stride
    4--32 passano da un blocco di attenzione (PAM lineare e CAM in parallelo) prima di essere
    sommate alle uscite del decoder. Il neck sul collo di bottiglia è l'identità oppure un
    quantizzatore vettoriale. Le frecce rosse sono le connessioni di salto.}
  \label{fig:maresunet}
  \vspace{1.2em}
  \resizebox{\linewidth}{!}{% Blocco di attenzione di MAResU-Net: conv 3x3, poi PAM lineare e CAM in parallelo,
% somma e proiezione 1x1. Entrambi i rami sono residuali con guadagno inizializzato a 0.
\begin{tikzpicture}[node distance=3.5mm and 6mm,
    b/.style={blocco, fill=white, draw=SapGrigio},
    somma/.style={circle, draw=SapRossoScuro, fill=white, inner sep=0pt, minimum size=4.2mm,
                  font=\footnotesize},
  ]
  \node[dato] (x) {$\vect{X}_l$\\{\scriptsize $h\times w\times C$}};
  \node[b, right=of x] (c3) {conv $3{\times}3$\\ReLU};
  \node[attn, above right=2mm and 7mm of c3, minimum width=2.6cm] (pam)
       {PAM lineare\\{\scriptsize $\vect{X}+\gamma\,\mathrm{LAM}(\vect{X})$}};
  \node[attn, below right=2mm and 7mm of c3, minimum width=2.6cm] (cam)
       {CAM\\{\scriptsize $\vect{X}+\beta\,\mathrm{CA}(\vect{X})$}};
  \node[b, right=of pam] (p1) {drop\\conv $1{\times}1$};
  \node[b, right=of cam] (c1) {drop\\conv $1{\times}1$};
  \node[somma, right=12mm of $(p1)!0.5!(c1)$] (s) {$+$};
  \node[b, right=of s] (o1) {drop\\conv $1{\times}1$};
  \node[dato, right=of o1] (y) {$\vect{S}_l$};
  \draw[flusso] (x) -- (c3);
  \draw[flusso] (c3.east) -- ++(3mm,0) |- (pam.west);
  \draw[flusso] (c3.east) -- ++(3mm,0) |- (cam.west);
  \draw[flusso] (pam) -- (p1);
  \draw[flusso] (cam) -- (c1);
  \draw[flusso] (p1.east) -| (s.north);
  \draw[flusso] (c1.east) -| (s.south);
  \draw[flusso] (s) -- (o1);
  \draw[flusso] (o1) -- (y);
  \node[nota, below=1mm of cam] {$\gamma=\beta=0$ all'inizializzazione};
\end{tikzpicture}
}
  \caption[Blocco di attenzione PAM-CAM]{Blocco di attenzione di \maresunet{}. Entrambi i rami
    sono residuali con guadagno iniziale nullo: all'inizio dell'addestramento il blocco calcola
    soltanto le convoluzioni.}
  \label{fig:pamcam}
\end{figure}

\begin{table}[t]
  \centering
  \small
  \caption[Costo dei moduli di attenzione]{Costo asintotico dei moduli di attenzione su una
    mappa con $N$ posizioni e $C$ canali ($d_k=C/8$ nel PAM).}
  \label{tab:costo-attenzione}
  \begin{tabular}{|l|c|c|l|}
    \hline
    \textbf{Modulo} & \textbf{Calcolo} & \textbf{Memoria aggiuntiva} & \textbf{Uso} \\
    \hline
    PAM dot-product & $\bigO{N^2(d_k+C)}$ & $\bigO{N^2}$ & solo mappe piccole \\
    PAM lineare (LAM) & $\bigO{N d_k C}$ & $\bigO{N(d_k+C)+d_kC}$ & tutti gli stadi \\
    CAM & $\bigO{NC^2}$ & $\bigO{C^2}$ & tutti gli stadi \\
    \hline
  \end{tabular}
\end{table}

% ---------------------------------------------------------------------
\section{Quantizzazione vettoriale}
\label{sec:vq}

\subsection{Codebook e assegnazione}

Un quantizzatore vettoriale~\cite{vandenoord2017vqvae} mantiene un \emph{codebook}
$\mathcal{E}=\{\vect{e}_1,\dots,\vect{e}_K\}\subset\R^{D}$ e sostituisce ogni vettore
$\vect{z}_e(u)\in\R^{D}$ prodotto dall'encoder con il codice più vicino:
\begin{equation}
  k^\star(u)=\argmin_{k}\norm{\vect{z}_e(u)-\vect{e}_k}_2^2,
  \qquad
  \vect{z}_q(u)=\vect{e}_{k^\star(u)} .
  \label{eq:vq-assegnazione}
\end{equation}
Lo spazio delle caratteristiche viene così partizionato nelle celle di Voronoi dei codici
(\cref{fig:vq}b), e il decoder riceve uno dei soli $K$ valori possibili per ogni posizione. La
distanza si calcola senza materializzare differenze, come
$\norm{\vect{z}}^2-2\vect{z}\T\vect{e}_k+\norm{\vect{e}_k}^2$.

\begin{figure}[t]
  \centering
  \resizebox{\linewidth}{!}{% Quantizzazione vettoriale: (a) flusso in avanti e gradiente straight-through;
% (b) geometria: assegnazione al codice più vicino, commitment e aggiornamento EMA.
\begin{tikzpicture}[node distance=4mm and 7mm]
  % ---------------- (a) flusso -------------------------------------------------
  \begin{scope}[local bounding box=A]
    \node[enc] (enc) {Encoder\\$f_\theta$};
    \node[dato, right=of enc] (ze) {$\vect{z}_e(u)\in\R^{D}$};
    \node[quant, right=of ze] (nn) {$k^\star=\argmin_k \norm{\vect{z}_e-\vect{e}_k}_2$};
    \node[dato, right=of nn] (zq) {$\vect{z}_q(u)=\vect{e}_{k^\star}$};
    \node[dec, right=of zq] (dec) {Decoder\\$g_\psi$};
    \node[blocco, fill=SapRossoVelo, below=6mm of nn, minimum width=3.6cm] (cb)
         {codebook $\{\vect{e}_1,\dots,\vect{e}_K\}\subset\R^{D}$};
    \draw[flusso] (enc) -- (ze);
    \draw[flusso] (ze) -- (nn);
    \draw[flusso] (nn) -- (zq);
    \draw[flusso] (zq) -- (dec);
    \draw[flusso] (cb) -- (nn);
    % gradiente straight-through
    \draw[grad] (dec.north) to[bend right=22]
      node[midway, above, font=\scriptsize, text=SapRosso]
      {$\nabla_{\vect{z}_q}\mathcal{L}$ copiato su $\vect{z}_e$ (straight-through)} (ze.north);
    \node[nota, anchor=north west] at ($(enc.south west)+(0,-3mm)$)
      {\textcolor{ColFreccia}{$\longrightarrow$} in avanti \quad
       \textcolor{SapRosso}{$\dashrightarrow$} gradiente};
  \end{scope}
  \node[font=\small\bfseries, anchor=south west] at ($(A.north west)+(0,1mm)$) {(a)};

  % ---------------- (b) geometria ----------------------------------------------
  \begin{scope}[shift={($(A.south west)+(1.0cm,-5.2cm)$)}, x=1.15cm, y=1.15cm, local bounding box=B]
    % cella del codice e_2: vettori assegnati e loro centroide
    \fill[SapRossoVelo] (2,1.6) circle (1.15);
    \foreach \q in {(1.35,1.15),(2.55,1.05),(1.6,2.35),(2.65,2.15),(1.2,1.85)}{
      \draw[SapGrigio, line width=0.5pt] \q ++(-2pt,-2pt) -- ++(4pt,4pt) \q ++(-2pt,2pt) -- ++(4pt,-4pt);
    }
    \coordinate (c) at (1.87,1.71);
    \draw[SapGrigio, very thick] (c) circle (2.4pt);
    % codici
    \foreach \p/\n/\a in {{(0,0)}/1/below, {(2,1.6)}/2/below right, {(4.6,0.2)}/3/below,
                          {(4.0,3.3)}/4/above, {(-0.3,2.9)}/5/above, {(6.6,3.1)}/6/above}{
      \fill[SapRossoScuro] \p circle (1.8pt) node[\a, font=\scriptsize] {$\vect{e}_{\n}$};
    }
    % vettore dell'encoder e commitment
    \coordinate (z) at (3.55,0.55);
    \draw[black, line width=0.7pt] (z) ++(-2.5pt,-2.5pt) -- ++(5pt,5pt) (z) ++(-2.5pt,2.5pt) -- ++(5pt,-5pt);
    \node[font=\scriptsize, anchor=north west] at (3.6,0.5) {$\vect{z}_e$};
    \draw[->, SapRosso, thick] (z) -- ($(z)!0.5!(2,1.6)$);
    \node[font=\scriptsize, text=SapRosso, anchor=south west] at (3.05,1.0) {commitment ($\beta$)};
    % aggiornamento EMA del codice
    \draw[->, SapRossoScuro, thick] (2,1.6) -- ($(2,1.6)!0.85!(c)$);
    % didascalie interne
    \node[font=\scriptsize, text=SapGrigio, align=center, anchor=south] at (2,2.8)
      {vettori assegnati a $\vect{e}_2$ ($\times$)\\e loro centroide ($\circ$)};
    \node[font=\scriptsize, text=SapRossoScuro, align=left, anchor=west] at (4.75,2.15)
      {EMA: $\vect{e}_2$ si sposta verso\\il centroide della sua cella};
    \draw[SapRossoScuro, densely dotted] (4.72,2.15) -- (2.12,1.66);
  \end{scope}
  \node[font=\small\bfseries, anchor=south west] at ($(B.north west)+(-1.0cm,1mm)$) {(b)};
\end{tikzpicture}
}
  \caption[Quantizzazione vettoriale]{Quantizzazione vettoriale. (a)~Flusso in avanti e
    gradiente \emph{straight-through}: il gradiente rispetto a $\vect{z}_q$ viene copiato su
    $\vect{z}_e$. (b)~Geometria nello spazio delle caratteristiche: il vettore $\vect{z}_e$ è
    assegnato al codice più vicino, il termine di commitment lo avvicina al codice e
    l'aggiornamento EMA sposta il codice verso il centroide dei vettori assegnati.}
  \label{fig:vq}
\end{figure}

\subsection{Gradiente straight-through e perdita}

L'$\argmin$ ha derivata nulla quasi ovunque. Lo stimatore
\emph{straight-through}~\cite{bengio2013straight} usa nel passo in avanti il valore quantizzato
e nel passo all'indietro l'identità:
\begin{equation}
  \tilde{\vect{z}}=\vect{z}_e+\sg\!\left[\vect{z}_q-\vect{z}_e\right],
  \qquad
  \tilde{\vect{z}}=\vect{z}_q,\quad
  \frac{\partial\tilde{\vect{z}}}{\partial\vect{z}_e}=I,
  \label{eq:straight-through}
\end{equation}
dove $\sg[\cdot]$ è l'operatore \emph{stop-gradient} (identità in avanti, derivata nulla). Il
codebook non riceve gradiente dal decoder e va appreso con un termine apposito. Si aggiunge
inoltre un termine di \emph{commitment} che impedisce alle uscite dell'encoder di allontanarsi
indefinitamente dai codici:
\begin{equation}
  \mathcal{L}_{\mathrm{VQ}}=
  \underbrace{\norm{\sg[\vect{z}_e]-\vect{z}_q}_2^2}_{\text{codebook}}
  +\beta\,\underbrace{\norm{\vect{z}_e-\sg[\vect{z}_q]}_2^2}_{\text{commitment}},
  \qquad
  \mathcal{L}=\mathcal{L}_{\mathrm{CE}}+\lambda\,\mathcal{L}_{\mathrm{VQ}},
  \label{eq:vq-loss}
\end{equation}
con le norme mediate sulle posizioni e sulle componenti, $\beta=0{,}25$ e $\lambda$ peso della
perdita ausiliaria.

\subsection{Aggiornamento EMA come K-means online}

In alternativa al termine \enquote{codebook} della~\eqref{eq:vq-loss}, i codici possono essere
aggiornati con medie mobili esponenziali~\cite{vandenoord2017vqvae,razavi2019vqvae2}. Sia, per
un batch, $n_k=\sum_u\Ind[k^\star(u)=k]$ il numero di vettori assegnati al codice $k$ e
$\vect{s}_k=\sum_u\Ind[k^\star(u)=k]\,\vect{z}_e(u)$ la loro somma. Con decadimento
$\gamma\in(0,1)$:
\begin{equation}
  N_k\leftarrow\gamma N_k+(1-\gamma)\,n_k,\qquad
  \vect{M}_k\leftarrow\gamma\vect{M}_k+(1-\gamma)\,\vect{s}_k,\qquad
  \vect{e}_k\leftarrow\vect{M}_k/\tilde N_k ,
  \label{eq:vq-ema}
\end{equation}
dove $\tilde N_k=\frac{N_k+\varepsilon}{\sum_j N_j+K\varepsilon}\sum_jN_j$ è una versione
regolarizzata (smoothing di Laplace) che evita divisioni per zero. Con l'EMA la perdita si
riduce al solo commitment, $\mathcal{L}_{\mathrm{VQ}}=\beta\norm{\vect{z}_e-\sg[\vect{z}_q]}^2$.

\begin{proposizione}[EMA e K-means]
\label{prop:ema-kmeans}
Trascurando lo smoothing e l'inizializzazione, dopo $t$ passi l'aggiornamento~\eqref{eq:vq-ema}
restituisce
\[
  \vect{e}_k^{(t)}=\frac{\sum_{\tau\le t}\gamma^{t-\tau}\vect{s}_k^{(\tau)}}
                        {\sum_{\tau\le t}\gamma^{t-\tau}n_k^{(\tau)}}
  =\argmin_{\vect{e}}\;\sum_{\tau\le t}\gamma^{t-\tau}\sum_{u}\Ind\!\left[k^{\star(\tau)}(u)=k\right]
   \norm{\vect{z}_e^{(\tau)}(u)-\vect{e}}_2^2 ,
\]
cioè il centroide dei vettori assegnati al codice $k$, pesati con uno sconto esponenziale
sull'età del batch. È un passo di Lloyd~\cite{lloyd1982kmeans} eseguito online.
\end{proposizione}
\begin{proof}
Srotolando le ricorrenze, $\vect{M}_k^{(t)}=(1-\gamma)\sum_{\tau}\gamma^{t-\tau}\vect{s}_k^{(\tau)}$ e
$N_k^{(t)}=(1-\gamma)\sum_\tau\gamma^{t-\tau}n_k^{(\tau)}$. Il fattore $(1-\gamma)$ si semplifica
nel rapporto. La funzione a destra è quadratica e convessa in $\vect{e}$, e annullandone il
gradiente,
$-2\sum_\tau\gamma^{t-\tau}\bigl(\vect{s}_k^{(\tau)}-n_k^{(\tau)}\vect{e}\bigr)=\vect{0}$, si
ottiene lo stesso rapporto.
\end{proof}

La \cref{prop:ema-kmeans} spiega perché lo stato EMA va trattato come stato persistente e non
come parametro ricreato a ogni passo, come accadeva nel 2022. Spiega anche il collasso tipico del
codebook: un codice che non riceve assegnazioni non si muove più, perché ha $n_k=0$ in ogni
batch, e diventa un \emph{codice morto}. Le contromisure adottate sono due
(\cref{alg:vq-ema}): l'\textbf{inizializzazione dai dati}, che pone i codici su vettori estratti
dal primo batch invece che in un intorno dell'origine lontano dalle attivazioni post-ReLU, e il
\textbf{riavvio dei codici morti}~\cite{lancucki2020robust}, che riposiziona su un vettore del
batch ogni codice con $N_k$ sotto una soglia.

\begin{algorithm}[t]
  \caption{Passo di addestramento del quantizzatore con EMA e riavvio dei codici morti}
  \label{alg:vq-ema}
  \begin{algorithmic}[1]
    \Require uscite dell'encoder $Z=\{\vect{z}_e(u)\}_{u}$, codebook $\{\vect{e}_k\}$, stato $\{N_k,\vect{M}_k\}$, decadimento $\gamma$, soglia $\tau_{\mathrm{morto}}$
    \If{primo passo}
      \State $\vect{e}_k\gets$ campione casuale di $Z$ (con piccolo rumore), $\vect{M}_k\gets\vect{e}_k$, $N_k\gets1$ per ogni $k$
    \EndIf
    \State $k^\star(u)\gets\argmin_k\norm{\vect{z}_e(u)-\vect{e}_k}^2$; \quad $\vect{z}_q(u)\gets\vect{e}_{k^\star(u)}$
    \State $\mathcal{L}_{\mathrm{VQ}}\gets\beta\cdot\mathrm{media}_u\norm{\vect{z}_e(u)-\sg[\vect{z}_q(u)]}^2$
    \State $N_k\gets\gamma N_k+(1-\gamma)n_k$; \quad $\vect{M}_k\gets\gamma\vect{M}_k+(1-\gamma)\vect{s}_k$
    \ForEach{$k$ con $N_k<\tau_{\mathrm{morto}}$}
      \State $\vect{M}_k\gets$ campione casuale di $Z$; \quad $N_k\gets1$ \Comment{riavvio}
    \EndFor
    \State $\vect{e}_k\gets\vect{M}_k/\tilde N_k$
    \State \Return $\vect{z}_e+\sg[\vect{z}_q-\vect{z}_e]$, $\mathcal{L}_{\mathrm{VQ}}$
  \end{algorithmic}
\end{algorithm}

\subsection{Diagnostica: perplexity e uso del codebook}

Sia $\bar p_k=n_k/\sum_jn_j$ la frequenza di assegnazione del codice $k$ in un batch. Si
definiscono
\begin{equation}
  \PPL=\exp\Bigl(-\sum_{k}\bar p_k\log\bar p_k\Bigr)\in[1,K],
  \qquad
  \mathrm{uso}=\frac{\abs{\{k:\bar p_k>0\}}}{K}\in\left[\tfrac1K,1\right].
  \label{eq:perplexity}
\end{equation}
La perplexity è il numero \emph{effettivo} di codici usati: vale 1 se tutti i vettori finiscono
sullo stesso codice e $K$ se l'uso è uniforme. Un codebook \enquote{sano}, nel senso della
\RQ{2}, ha perplexity alta e uso vicino a 1 durante tutto l'addestramento.

\subsection{Capacità del collo di bottiglia e aggiramento}
\label{sec:vq-aggiramento}

La quantizzazione limita l'informazione che attraversa il collo di bottiglia. Con $h\times w$
posizioni e $K$ codici,
\begin{equation}
  I(\vect{x};\vect{z}_q)\;\le\;\Hent(\vect{z}_q)\;\le\;hw\log_2K\ \text{bit}.
  \label{eq:capacita}
\end{equation}
Per un ritaglio di $512\times512$ pixel il collo di bottiglia di \maresunet{} è a stride 32
($16\times16$ posizioni). Con $K=512$ porta al più $256\cdot9=\num{2304}$ bit, cioè 288 byte per
ritaglio. La sola connessione di salto a stride 4 trasporta invece
$128\cdot128\cdot64\approx10^6$ valori reali. Il collo di bottiglia vincola quindi solo il
percorso più stretto della rete, che non è l'unico. Ne segue il risultato seguente.

\begin{proposizione}[Aggiramento del collo di bottiglia]
\label{prop:aggiramento}
Si consideri la~\eqref{eq:maresunet} con il quantizzatore solo sul neck e si supponga che la
perdita di segmentazione $\mathcal{L}_{\mathrm{CE}}$ ammetta un minimo globale con un decoder
che non dipende da $\vect{B}$. Allora la perdita totale~\eqref{eq:vq-loss} ammette un minimo
globale in cui il codebook usa un solo codice ($\PPL=1$).
\end{proposizione}
\begin{proof}
Si parta dal minimo di $\mathcal{L}_{\mathrm{CE}}$ con decoder indipendente da $\vect{B}$. Le
mappe $\vect{S}_1,\vect{S}_2,\vect{S}_3$ non dipendono dallo stadio $\vect{E}_4$. Si possono
quindi modificare i parametri dell'ultimo stadio (per esempio azzerando il guadagno dell'ultima
normalizzazione) in modo che $\vect{E}_4$, e quindi $\vect{z}_e$, sia costante e pari a
$\vect{c}$, senza cambiare $\mathcal{L}_{\mathrm{CE}}$. Ponendo $\vect{e}_1=\vect{c}$, il
commitment è nullo e tutti i vettori sono assegnati a $\vect{e}_1$. Entrambi i termini, che sono
non negativi, sono al loro minimo.
\end{proof}

La proposizione non afferma che l'ottimizzazione raggiunga quella soluzione. Afferma che
\emph{nulla nell'obiettivo la penalizza}: con le sole connessioni di salto libere, un risultato
negativo della \maresunetvq{} non dice nulla sull'utilità della quantizzazione. Per questo la
\RQ{2} misura la perplexity durante l'addestramento e prevede una variante in cui anche le
connessioni di salto sono quantizzate, così che ogni percorso dall'encoder al decoder passi da
un codebook e l'ipotesi della \cref{prop:aggiramento} cada.

% ---------------------------------------------------------------------
\section{Auto-distillazione senza etichette}
\label{sec:dino}

\subsection{Student, teacher e distribuzioni}

DINO~\cite{caron2021dino} addestra uno \emph{student} $g_s=h_s\circ\GAP\circ f_s$ a riprodurre
l'uscita di un \emph{teacher} $g_t=h_t\circ\GAP\circ f_t$ con la stessa architettura. Qui $f$ è
l'encoder (lo stesso che verrà usato in \maresunet{}), $\GAP$ la media globale sulle posizioni e
$h$ una testa di proiezione: un MLP a tre strati (2048, 2048, 256 unità con GELU),
normalizzazione $\ell_2$ e uno strato lineare a peso normalizzato verso $K=\num{16384}$
\enquote{prototipi}. Le uscite diventano distribuzioni con due temperature:
\begin{equation}
  P_s(\vect{x})=\softmax\!\left(\frac{g_s(\vect{x})}{\tau_s}\right),\qquad
  P_t(\vect{x})=\softmax\!\left(\frac{g_t(\vect{x})-\vect{c}}{\tau_t}\right),
  \label{eq:dino-distribuzioni}
\end{equation}
dove $\vect{c}\in\R^K$ è il \emph{centro} e $\tau_t<\tau_s$ rende la distribuzione del teacher
più concentrata (\emph{sharpening}).

\subsection{Multi-crop e perdita}

Da ogni immagine si estraggono due viste globali $\vect{x}^{g_1},\vect{x}^{g_2}$ ($224\times224$,
scala tra il 14\% e il 100\% dell'area) e $M$ viste locali $\vect{x}^{l_1},\dots,\vect{x}^{l_M}$
($96\times96$, scala tra il 5\% e il 40\%), con trasformazioni cromatiche e geometriche
casuali. Il teacher vede solo le viste globali, lo student tutte. Per un'immagine la perdita è
\begin{equation}
  \ell(\vect{x})=\frac{1}{\abs{\mathcal{P}}}\sum_{(a,b)\in\mathcal{P}}
  \Hent\bigl(P_t(\vect{x}^{a}),P_s(\vect{x}^{b})\bigr),
  \qquad
  \mathcal{P}=\bigl\{(a,b):a\in\{g_1,g_2\},\ b\in V,\ b\neq a\bigr\},
  \label{eq:dino-loss}
\end{equation}
con $V$ l'insieme di tutte le viste, $\Hent(p,q)=-\sum_kp_k\log q_k$ e
$\abs{\mathcal{P}}=2(M+1)$. Sul batch si prende la media
$\mathcal{L}=\frac1B\sum_i\ell(\vect{x}_i)$. La perdita incoraggia corrispondenze
\enquote{da locale a globale}: lo student deve dedurre da una porzione dell'immagine la
distribuzione che il teacher assegna a una vista ampia.

\begin{osservazione}
La normalizzazione~\eqref{eq:dino-distribuzioni} è \emph{per campione}: ogni immagine ha la sua
distribuzione su $K$ prototipi. Calcolare la softmax sull'intero batch appiattito, come avveniva
nel codice del 2022, produce un'unica distribuzione su $B\cdot K$ elementi (lì addirittura sui
pixel dell'uscita densa del decoder). Le immagini del batch competono allora per la stessa massa
di probabilità e l'obiettivo non è più quello di DINO.
\end{osservazione}

\begin{figure}[t]
  \centering
  \resizebox{\linewidth}{!}{% Auto-distillazione DINO con multi-crop: teacher sulle viste globali, student su tutte;
% centering e sharpening sul teacher, stop-gradient, teacher come EMA dello student.
\begin{tikzpicture}[node distance=4mm and 6mm,
    vista/.style={dato, minimum width=1.2cm, minimum height=6mm, font=\scriptsize},
  ]
  \node[dato, minimum height=12mm] (x) {immagine\\$\vect{x}$};

  \node[vista, above right=3mm and 9mm of x] (g) {$\vect{x}^{g}_{1}, \vect{x}^{g}_{2}$\\$224^2$};
  \node[vista, below right=3mm and 9mm of x] (l) {$\vect{x}^{l}_{1..M}$\\$96^2$};
  \draw[flusso] (x.east) -- ++(3mm,0) |- (g.west);
  \draw[flusso] (x.east) -- ++(3mm,0) |- (l.west);
  \node[nota, above=0mm of g] {viste globali};
  \node[nota, below=0mm of l] {viste locali};

  % teacher
  \node[enc, right=14mm of g, minimum width=2.5cm, fill=SapGrigioChiaro, draw=SapGrigio] (t)
       {teacher $h_t\circ f_t$\\{\scriptsize $\theta_t$, nessun gradiente}};
  \node[blocco, fill=white, right=of t] (tc) {$\softmax\!\big((\cdot-\vect{c})/\tau_t\big)$};
  \node[dato, right=of tc] (pt) {$P_t$};

  % student
  \node[enc, right=14mm of l, minimum width=2.5cm] (s) {student $h_s\circ f_s$\\{\scriptsize $\theta_s$}};
  \node[blocco, fill=white, right=of s] (sc) {$\softmax\!\big(\cdot/\tau_s\big)$};
  \node[dato, right=of sc] (ps) {$P_s$};

  \draw[flusso] (g) -- (t);
  \draw[flusso] (l) -- (s);
  \draw[flusso] (g.east) -- ++(4mm,0) |- ([yshift=1.5mm]s.west);
  \draw[flusso] (t) -- (tc);
  \draw[flusso] (tc) -- (pt);
  \draw[flusso] (s) -- (sc);
  \draw[flusso] (sc) -- (ps);

  % loss
  \node[quant, right=10mm of $(pt)!0.5!(ps)$, align=center] (loss)
       {$\displaystyle\sum_{g}\sum_{v\neq g}\Hent\big(P_t(\vect{x}^g),P_s(\vect{x}^v)\big)$};
  \draw[flusso] (pt) -| (loss.north);
  \draw[flusso] (ps) -| (loss.south);
  \draw[grad] (loss.south) ++(-4mm,0) |- ([yshift=-3mm]ps.south) -| (s.south)
       node[pos=0.25, below, font=\scriptsize, text=SapRosso] {retropropagazione solo nello student};
  \node[font=\scriptsize, text=SapRosso, right=0.5mm of pt, yshift=3mm] {sg};

  % EMA and centering
  \draw[->, SapRossoScuro, thick] (s.north) -- (t.south)
       node[midway, right, font=\scriptsize, align=left]
       {$\theta_t\leftarrow m\,\theta_t+(1-m)\,\theta_s$};
  \draw[->, SapGrigio, densely dotted, thick] (tc.north) to[out=110, in=70, looseness=4]
       node[above, font=\scriptsize, text=SapGrigio] {$\vect{c}\leftarrow\lambda\vect{c}+(1-\lambda)\,\overline{h_t f_t(\vect{x}^{g})}$} (tc.north east);
\end{tikzpicture}
}
  \caption[Auto-distillazione DINO]{Auto-distillazione con multi-crop. Il teacher vede solo le
    viste globali e non riceve gradiente (sg). La sua uscita è centrata e resa più concentrata
    con una temperatura bassa. Il teacher è una media mobile esponenziale dello student e il
    centro è una media mobile delle uscite del teacher.}
  \label{fig:dino}
\end{figure}

\subsection{Teacher come media mobile e centering}

Il teacher non viene ottimizzato: i suoi parametri inseguono quelli dello student con una media
mobile esponenziale, alla maniera di un \emph{mean teacher}~\cite{tarvainen2017meanteacher} o di
una media di Polyak~\cite{polyak1992averaging}:
\begin{equation}
  \theta_t\leftarrow m\,\theta_t+(1-m)\,\theta_s,\qquad
  m(\tau)=1-(1-m_0)\,\frac{1+\cos(\pi\tau/T)}{2},
  \label{eq:dino-ema}
\end{equation}
con $m_0=0{,}996$ che cresce fino a 1 alla fine dell'addestramento (passo $\tau$ su $T$). Il
teacher nasce come copia esatta dello student, e la media include le statistiche della batch
normalization. Il centro è a sua volta una media mobile delle uscite del teacher sul batch:
\begin{equation}
  \vect{c}\leftarrow\lambda\,\vect{c}+(1-\lambda)\,\frac{1}{2B}\sum_{i=1}^{B}\sum_{a\in\{g_1,g_2\}}g_t(\vect{x}_i^{a}),
  \qquad\lambda=0{,}9 .
  \label{eq:dino-centro}
\end{equation}

\subsection{Collasso e sua diagnosi}
\label{sec:collasso}

L'entropia incrociata si scompone come
$\Hent(P_t,P_s)=\Hent(P_t)+\KL(P_t\,\Vert\,P_s)$. Se il teacher produce la stessa distribuzione
per ogni ingresso, lo student la può imitare perfettamente ($\KL=0$) senza aver imparato nulla:
è il \emph{collasso}. Ne esistono due forme:
\begin{itemize}
  \item \textbf{collasso su una dimensione}: $P_t$ è quasi un vettore indicatore, uguale per
  tutte le immagini, e $\Hent(P_t)\to0$;
  \item \textbf{collasso uniforme}: $P_t$ è quasi uniforme per tutte le immagini, e
  $\Hent(P_t)\to\log K$ ($\approx9{,}70$ nat per $K=\num{16384}$).
\end{itemize}
Il \emph{centering}~\eqref{eq:dino-centro} sottrae la componente comune a tutte le immagini e
impedisce a una dimensione di dominare, ma da solo spinge verso l'uniforme. Lo
\emph{sharpening} ($\tau_t$ piccola) ha l'effetto opposto. DINO bilancia le due forze.

Entrambe le forme di collasso hanno in comune l'indipendenza dell'uscita dall'ingresso, che si
misura con l'informazione mutua tra immagine e prototipo:
\begin{equation}
  I=\Hent\bigl(\E_{\vect{x}}[P_t(\vect{x})]\bigr)-\E_{\vect{x}}\bigl[\Hent(P_t(\vect{x}))\bigr]\;\ge\;0 .
  \label{eq:dino-mi}
\end{equation}
Il collasso, in entrambe le forme, corrisponde a $I\to0$. Nel monitoraggio si registra a ogni
epoca l'entropia media del teacher $\E_{\vect{x}}[\Hent(P_t)]$, che tende a $0$ o a $\log K$ in
caso di collasso.
\Segnaposto{Aggiungere al monitoraggio anche l'entropia della distribuzione media
$\Hent(\E_{\vect{x}}[P_t])$, così da stimare la~\eqref{eq:dino-mi}.}

\subsection{Schedule}

Il learning rate cresce linearmente per 10 epoche fino a $\eta_0=5\cdot10^{-4}\cdot B/256$
(regola di scala lineare) e poi decresce con un coseno fino a $10^{-6}$. Il \emph{weight decay}
cresce con un coseno da 0,04 a 0,4. La temperatura del teacher cresce linearmente da 0,04 a 0,07
nelle prime 30 epoche, mentre $\tau_s=0{,}1$ resta fissa. Uno schedule coseno da $v_0$ a
$v_{\mathrm{fin}}$ su $T$ passi è
$v(\tau)=v_{\mathrm{fin}}+(v_0-v_{\mathrm{fin}})\frac{1+\cos(\pi\tau/T)}{2}$. Il
\cref{alg:dino} riassume un passo di addestramento.

\begin{algorithm}[t]
  \caption{Un passo di addestramento DINO}
  \label{alg:dino}
  \begin{algorithmic}[1]
    \Require batch $\{\vect{x}_i\}_{i=1}^{B}$, student $\theta_s$, teacher $\theta_t$, centro $\vect{c}$, passo $\tau$
    \State $(\eta,\mathrm{wd},m,\tau_t)\gets$ schedule al passo $\tau$
    \ForEach{immagine $\vect{x}_i$}
      \State genera $\vect{x}_i^{g_1},\vect{x}_i^{g_2}$ ($224^2$) e $\vect{x}_i^{l_1},\dots,\vect{x}_i^{l_M}$ ($96^2$)
    \EndFor
    \State $P_t\gets\softmax\bigl((g_t(\vect{x}^{g})-\vect{c})/\tau_t\bigr)$ \Comment{senza gradiente}
    \State $P_s\gets\softmax\bigl(g_s(\vect{x}^{v})/\tau_s\bigr)$ per tutte le viste $v$
    \State $\mathcal{L}\gets\frac1B\sum_i\frac{1}{\abs{\mathcal{P}}}\sum_{(a,b)\in\mathcal{P}}\Hent\bigl(P_t(\vect{x}_i^{a}),P_s(\vect{x}_i^{b})\bigr)$
    \State aggiorna $\theta_s$ con AdamW($\eta$, wd) e clipping della norma del gradiente
    \State $\theta_t\gets m\,\theta_t+(1-m)\,\theta_s$
    \State $\vect{c}\gets\lambda\vect{c}+(1-\lambda)\,\mathrm{media}_{i,a}\,g_t(\vect{x}_i^{a})$
  \end{algorithmic}
\end{algorithm}

\subsection{Augmentation e immagini aeree}
\label{sec:augmentation-teoria}

Le augmentation definiscono le invarianze che la rappresentazione impara: tutto ciò che varia
tra due viste della stessa immagine viene trattato come rumore. Nelle immagini naturali le
trasformazioni cromatiche forti (variazioni di luminosità, contrasto e saturazione, scala di
grigi, solarizzazione) insegnano a ignorare l'illuminazione. Nelle immagini aeree, però, il
colore è spesso l'indizio principale: in falso colore IR-R-G la vegetazione è rossa per la
riflettanza nel vicino infrarosso. Rendere la rete invariante al colore può quindi eliminare
informazione utile. D'altra parte, le immagini aeree non hanno un'orientazione canonica, e le
rotazioni di \ang{90} e le riflessioni (il gruppo diedrale $D_4$) sono invarianze corrette che le
augmentation di DINO non usano. La variante \enquote{aerea} studiata nella \RQ{4} attenua le
trasformazioni cromatiche e aggiunge le trasformazioni di $D_4$.

\subsection{Dal pre-addestramento alla segmentazione}

Al termine del pre-addestramento si conserva solo l'encoder dello student, $f_s$, che inizializza
l'encoder di \maresunet{}. Attenzione, decoder e testa partono da un'inizializzazione casuale.

\begin{osservazione}[Obiettivo globale, compito denso]
\label{oss:globale-denso}
La perdita~\eqref{eq:dino-loss} agisce dopo la media globale $\GAP$: premia rappresentazioni
che descrivono l'immagine nel suo insieme, non le singole posizioni. Il compito a valle è invece
denso. Nulla nell'obiettivo garantisce che le mappe intermedie $\vect{E}_1,\vect{E}_2$, che
alimentano le connessioni di salto, siano informative a livello di pixel. È una possibile
ragione di un divario tra DINO e ImageNet indipendente dalla quantità di dati, e motiva la
variante densa discussa nel \cref{cap:prospettive}.
\end{osservazione}

% ---------------------------------------------------------------------
\section{Fusione dell'informazione altimetrica}
\label{sec:ndsm}

\subsection{Dal DSM all'nDSM}

Il modello digitale di superficie (DSM) dà la quota assoluta della superficie visibile, quindi
anche di tetti e chiome. Il modello normalizzato (nDSM) sottrae la quota del terreno (DTM) e
restituisce l'altezza sul terreno, che è l'informazione utile a distinguere un tetto da una
strada o un albero da un prato. Quando l'nDSM non è disponibile, il DTM si può stimare con
un'\emph{apertura morfologica}~\cite{soille2003morphology}. Detti
$\varepsilon_B(f)(u)=\min_{v\in B}f(u+v)$ ed $\delta_B(f)(u)=\max_{v\in B}f(u-v)$ l'erosione e
la dilatazione con elemento strutturante $B$, si ha
\begin{equation}
  \widehat{\mathrm{DTM}}=\gamma_B(\mathrm{DSM})=\delta_B\bigl(\varepsilon_B(\mathrm{DSM})\bigr),
  \qquad
  \mathrm{nDSM}=\mathrm{DSM}-\widehat{\mathrm{DTM}} .
  \label{eq:apertura}
\end{equation}
L'apertura elimina le strutture in rilievo più strette di $B$ (edifici, alberi, auto) e conserva
l'andamento del terreno. Essendo \emph{anti-estensiva} ($\gamma_B(f)\le f$), garantisce
$\mathrm{nDSM}\ge0$ per costruzione. In pratica si calcola il minimo del DSM su celle di circa
\SI{1}{\meter}, si applica l'apertura con una finestra di circa \SI{30}{\meter}, si smussa il
risultato e lo si riporta alla risoluzione originale. Il limite noto è che gli edifici più larghi
della finestra lasciano parte della loro altezza nel terreno stimato. L'nDSM viene codificato in
decimetri su 8 bit, $h=\min(\mathrm{round}(10\cdot\mathrm{nDSM}),255)$, con saturazione a
\SI{25,5}{\meter}.

\subsection{Fusione precoce ed estensione dei pesi}

Nella fusione precoce l'nDSM normalizzato con media e deviazione standard del training,
$\tilde h=(h-\mu_h)/\sigma_h$, diventa un quarto canale di ingresso. Il primo strato di un
encoder pre-addestrato su tre canali, con pesi $W\in\R^{k\times k\times3\times C}$, va esteso a
$W'=[\,W\mid\vect{w}_4\,]\in\R^{k\times k\times4\times C}$.

\begin{proposizione}[Estensione con pesi nulli]
\label{prop:inflazione}
Se $\vect{w}_4=\vect{0}$, la rete con ingresso $[\vect{x};\tilde h]$ calcola esattamente la
stessa funzione della rete RGB con ingresso $\vect{x}$, per ogni $\tilde h$. Il gradiente
rispetto a $\vect{w}_4$ è però in generale non nullo, quindi l'addestramento può imparare a
usare l'altezza partendo dal modello RGB.
\end{proposizione}
\begin{proof}
Per linearità della convoluzione,
$W'*[\vect{x};\tilde h]=W*\vect{x}+\vect{w}_4*\tilde h=W*\vect{x}$ quando $\vect{w}_4=\vect{0}$.
Detto $\vect{\delta}$ il gradiente della perdita rispetto all'uscita del primo strato,
$\partial\mathcal{L}/\partial\vect{w}_4[a,b,:]=\sum_u\tilde h(u+(a,b))\,\vect{\delta}(u)$, che è
non nullo appena l'altezza è correlata con l'errore.
\end{proof}

L'alternativa $\vect{w}_4=\frac13\sum_cW_{:,:,c,:}$ tratta l'altezza come un canale di luminanza
aggiuntivo. Perturba la funzione iniziale, ma dà subito al canale un peso dell'ordine di quelli
RGB. Il confronto tra le due inizializzazioni fa parte della \RQ{6}.

% ---------------------------------------------------------------------
\section{Metriche di valutazione}
\label{sec:metriche}

\subsection{Matrice di confusione e metriche per classe}

Tutte le metriche derivano dalla matrice di confusione $M\in\N^{K\times K}$, accumulata su
\emph{tutti} i pixel etichettati dello split:
$M_{ij}=\abs{\{u\in\Omega_v : y_u=i,\ \hat y_u=j\}}$. Per la classe $k$:
\begin{gather}
  \mathrm{TP}_k=M_{kk},\qquad
  \mathrm{FP}_k=\sum_{i\neq k}M_{ik},\qquad
  \mathrm{FN}_k=\sum_{j\neq k}M_{kj},\nonumber\\
  P_k=\frac{\mathrm{TP}_k}{\mathrm{TP}_k+\mathrm{FP}_k},\qquad
  R_k=\frac{\mathrm{TP}_k}{\mathrm{TP}_k+\mathrm{FN}_k},\nonumber\\
  F1_k=\frac{2\,\mathrm{TP}_k}{2\,\mathrm{TP}_k+\mathrm{FP}_k+\mathrm{FN}_k},\qquad
  \mathrm{IoU}_k=\frac{\mathrm{TP}_k}{\mathrm{TP}_k+\mathrm{FP}_k+\mathrm{FN}_k},
  \label{eq:metriche}
\end{gather}
e l'accuratezza complessiva è $\mathrm{OA}=\operatorname{tr}(M)/\sum_{ij}M_{ij}$. Precision e
recall corrispondono a \emph{correctness} e \emph{completeness} nella terminologia ISPRS.
L'IoU (indice di Jaccard~\cite{everingham2010pascal}) e l'F1 sono legati da una relazione
monotona:
\begin{equation}
  \mathrm{IoU}_k=\frac{F1_k}{2-F1_k},
\end{equation}
che si verifica sostituendo $T=\mathrm{TP}_k$ ed $E=\mathrm{FP}_k+\mathrm{FN}_k$ nelle
definizioni. Le due metriche ordinano quindi i modelli allo stesso modo classe per classe, ma
non necessariamente in media. mIoU e mF1 sono medie sulle classi presenti nel riferimento.
Vengono riportate con tutte le classi e senza clutter, che in Vaihingen è quasi assente.

\subsection{Perché una matrice di confusione globale}

Mediare le metriche calcolate batch per batch non equivale a calcolarle sull'intero split,
perché la media di rapporti è diversa dal rapporto delle somme. L'esempio della
\cref{tab:media-batch} mostra lo scostamento su una classe rara. Usare \codice{nanmean} per
saltare i batch in cui la classe è assente peggiora le cose, perché fa pesare un batch con un
solo pixel della classe quanto un batch che ne contiene migliaia.

\begin{table}[t]
  \centering
  \small
  \caption[Media per batch e matrice di confusione globale]{IoU di una classe rara calcolato
    come media per batch e dalla matrice di confusione globale.}
  \label{tab:media-batch}
  \begin{tabular}{|l|r|r|r|c|}
    \hline
    & \textbf{TP} & \textbf{FP} & \textbf{FN} & \textbf{IoU} \\
    \hline
    Batch A (un'auto piccola) & 1 & 0 & 0 & 1,000 \\
    Batch B (parcheggio) & 90 & 60 & 50 & 0,450 \\
    \hline
    Media per batch & & & & 0,725 \\
    \rowcolor{RigaRiferimento}
    Matrice globale & 91 & 60 & 50 & 0,453 \\
    \hline
  \end{tabular}
\end{table}

\subsection{Etichette piene ed \emph{eroded}}

Il benchmark ISPRS fornisce anche etichette \emph{eroded}, in cui una fascia di pixel lungo i
contorni degli oggetti è marcata come non valida. Il motivo è che la posizione esatta di un
contorno è incerta anche per l'annotatore. Valutare sulle etichette eroded significa sostituire
$\Omega_v$ con $\Omega_v^{\mathrm{er}}\subset\Omega_v$. Il confronto tra le due versioni
quantifica quanta parte dell'errore si concentra sui bordi. Le definizioni coincidono con quelle
dello strumento di valutazione ufficiale ISPRS.

% ---------------------------------------------------------------------
\section{Inferenza su tile intere}
\label{sec:finestra}

Le tile di test (fino a $6000\times6000$ pixel) non entrano in memoria in un solo passo. Si
predicono quindi con finestre di lato $w$ e passo $s<w$, e le probabilità delle finestre
sovrapposte vengono fuse. Vicino ai bordi di una finestra il contesto è troncato e la
predizione è meno affidabile: le probabilità vengono quindi pesate con una finestra di Hann
separabile,
\begin{equation}
  \omega(u)=\omega_1(u_x)\,\omega_1(u_y),\qquad
  \omega_1(t)=\sin^2\!\left(\frac{\pi(t+1)}{w+1}\right),\quad t=0,\dots,w-1,
  \label{eq:hann}
\end{equation}
con un minimo di $10^{-3}$. Dette $o_j$ le origini delle finestre e $\vect{p}_j$ le loro
probabilità,
\begin{equation}
  P(u)=\frac{\sum_j\omega(u-o_j)\,\vect{p}_j(u-o_j)}{\sum_j\omega(u-o_j)},
  \qquad
  \hat y_u=\argmax_kP_k(u).
  \label{eq:fusione}
\end{equation}
Con passo $s=w/2$ e nel limite continuo vale $\sin^2(\pi t/w)+\sin^2(\pi(t+w/2)/w)=1$: le
finestre formano una partizione dell'unità e il denominatore della~\eqref{eq:fusione} è
costante lontano dai bordi della tile (\cref{fig:hann}). La normalizzazione serve allora solo
vicino ai bordi e dove l'ultima finestra di ogni riga e colonna è allineata al bordo della tile.
Le tile più piccole della finestra vengono estese per riflessione.

\begin{figure}[t]
  \centering
  % Pesi delle finestre di Hann con sovrapposizione del 50%: la somma vale 1 lontano dai bordi.
\begin{tikzpicture}
  \begin{axis}[
      width=0.9\linewidth, height=4.8cm,
      xlabel={posizione lungo la tile (in unità di lato della finestra $w$)},
      ylabel={peso},
      xmin=0, xmax=3, ymin=0, ymax=1.25,
      xtick={0,0.5,...,3}, ytick={0,0.5,1},
      grid=major, grid style={SapGrigioChiaro},
      legend style={font=\footnotesize, draw=SapGrigioTenue, at={(0.5,1.03)}, anchor=south},
      legend columns=3,
      tick label style={font=\footnotesize}, label style={font=\small},
      samples=200,
    ]
    \addplot[SapRossoTenue, thick, domain=0:1] {sin(deg(pi*x))^2};
    \addlegendentry{$\omega$ delle singole finestre}
    \addplot[SapRossoTenue, thick, domain=0.5:1.5, forget plot] {sin(deg(pi*(x-0.5)))^2};
    \addplot[SapRossoTenue, thick, domain=1:2, forget plot] {sin(deg(pi*(x-1)))^2};
    \addplot[SapRossoTenue, thick, domain=1.5:2.5, forget plot] {sin(deg(pi*(x-1.5)))^2};
    \addplot[SapRossoTenue, thick, domain=2:3, forget plot] {sin(deg(pi*(x-2)))^2};
    \addplot[SapRosso, very thick, domain=0.5:2.5] {1};
    \addlegendentry{somma dei pesi}
    \addplot[SapGrigio, very thick, dashed, domain=0:0.5] {sin(deg(pi*x))^2};
    \addlegendentry{bordo: compensato dalla normalizzazione}
    \addplot[SapGrigio, very thick, dashed, domain=2.5:3, forget plot] {sin(deg(pi*(x-2)))^2};
  \end{axis}
\end{tikzpicture}

  \caption[Fusione delle finestre con pesi di Hann]{Pesi di Hann con sovrapposizione del 50\%:
    la somma è costante lontano dai bordi, quindi la fusione~\eqref{eq:fusione} è una media
    convessa che privilegia il centro di ogni finestra.}
  \label{fig:hann}
\end{figure}

% =====================================================================
%  Capitolo 4 — Metodologia sperimentale
% =====================================================================
\chapter{Metodologia sperimentale}
\label{cap:metodologia}

Questo capitolo traduce le domande di ricerca in un disegno sperimentale: ipotesi, dati,
protocolli di split, configurazioni di addestramento e di valutazione, implementazione, budget
computazionale e analisi statistica. Il principio guida è che ogni confronto cambi \emph{un solo}
fattore alla volta e che ogni differenza venga riportata con la sua variabilità.

% ---------------------------------------------------------------------
\section{Ipotesi e disegno sperimentale}
\label{sec:disegno}

Le domande di ricerca della \cref{sec:domande} si traducono in ipotesi verificabili. Nel seguito
$\Delta$ indica la differenza di mIoU di test tra la variante e il suo riferimento, mediata sui
seed. Ogni ipotesi è formulata in modo che anche il suo rifiuto sia un risultato informativo.

\begin{ipotesi}[\RQ{1}]\label{ip:lam}
A parità di encoder e protocollo, le connessioni di salto con attenzione lineare migliorano la
mIoU rispetto alle connessioni identità: $\Delta_{\mathrm{LAM}}>0$.
\end{ipotesi}
\begin{ipotesi}[\RQ{2}]\label{ip:vq}
Un codebook sano (perplexity alta durante tutto l'addestramento) non peggiora la mIoU:
$\Delta_{\mathrm{VQ}}\ge0$. Con le sole connessioni di salto libere la perplexity diminuisce
durante l'addestramento (\cref{prop:aggiramento}).
\end{ipotesi}
\begin{ipotesi}[\RQ{3}]\label{ip:dino}
Il vantaggio relativo del pre-addestramento DINO su ImageNet cresce al diminuire della frazione
di etichette: $\Delta_{\mathrm{DINO}}(1\%)>\Delta_{\mathrm{DINO}}(100\%)$.
\end{ipotesi}
\begin{ipotesi}[\RQ{4}]\label{ip:aug}
La variante \enquote{aerea} delle augmentation produce un encoder migliore della variante
standard, soprattutto con poche etichette.
\end{ipotesi}
\begin{ipotesi}[\RQ{6}]\label{ip:ndsm}
L'nDSM come quarto canale migliora la mIoU, con i guadagni maggiori sulle classi edifici e alberi.
\end{ipotesi}

La \RQ{5} è descrittiva: misura costi e accuratezze senza un'ipotesi direzionale. La
\cref{tab:disegno} collega ogni domanda alla variabile manipolata, al riferimento e alla
campagna di esperimenti che la esegue.

\begin{table}[t]
  \centering
  \small
  \caption[Disegno sperimentale per domanda di ricerca]{Disegno sperimentale: per ogni domanda di
    ricerca, il fattore che varia, il riferimento, la campagna che la esegue e il numero di
    passi (pre-addestramenti, training, benchmark) con tre seed.}
  \label{tab:disegno}
  \begin{tabular}{|c|p{3.5cm}|p{2.7cm}|l|c|}
    \hline
    \textbf{RQ} & \textbf{Fattore} & \textbf{Riferimento} & \textbf{Campagna} & \textbf{Passi} \\
    \hline
    \RQ{1} & connessioni di salto: LAM / identità & \maresunet{} R18 ImageNet & \codice{rq1\_attention} & 6 \\
    \hline
    \RQ{2} & neck: identità / VQ; encoder R18/R34; pre-addestr. ImageNet/DINO & stessa configurazione senza VQ & \codice{thesis\_grid} & 26 \\
    \hline
    \RQ{3} & frazione di etichette $\{1, 10, 25, 100\}\%$ $\times$ pre-addestr. & ImageNet, stessa frazione & \codice{rq3\_low\_label} & 38 \\
    \hline
    \RQ{4} & augmentation DINO: standard / aerea & DINO standard & \codice{rq4\_aerial} & 14 \\
    \hline
    \RQ{5} & encoder R18/R34/R50, Swin-T, ViT & \maresunet{} R18 & \makecell[l]{\codice{rq5\_efficiency}\\\codice{transformers}} & 10 + 10 \\
    \hline
    \RQ{6} & ingresso RGB / RGB + nDSM (ufficiale, calcolato) & RGB & \makecell[l]{\codice{rq6\_ndsm}\\\codice{vaihingen}} & 9 + 9 \\
    \hline
    -- & protocollo del 2022 & tab.~\ref{tab:risultati-2022} & \codice{thesis2022} & 3 \\
    \hline
  \end{tabular}
\end{table}

% ---------------------------------------------------------------------
\section{Dati}
\label{sec:dati}

\subsection{ISPRS Potsdam e Vaihingen}

Il dataset principale è ISPRS Potsdam, il secondo benchmark è ISPRS
Vaihingen~\cite{rottensteiner2012isprs,isprs2dlabeling} (\cref{tab:dataset}). Le classi e la
palette delle maschere sono le stesse per i due dataset (\cref{tab:classi}). Rispetto al codice
del 2022, l'ordine degli indici segue la convenzione ISPRS/MMSegmentation, con il clutter per
ultimo: così il valore ignorato (255) non può coincidere per errore con una classe reale.

\begin{table}[t]
  \centering
  \small
  \caption[Caratteristiche dei dataset ISPRS]{Caratteristiche dei dataset ISPRS usati.}
  \label{tab:dataset}
  \begin{tabular}{|l|p{4.6cm}|p{4.6cm}|}
    \hline
    & \textbf{Potsdam} & \textbf{Vaihingen} \\
    \hline
    Aree & 38 tile $6000\times6000$ & 33 aree, circa $1900\times2500$ \\
    GSD & \SI{5}{\centi\meter} & \SI{9}{\centi\meter} \\
    Immagine & RGB (\codice{2\_Ortho\_RGB}) & IR-R-G in falso colore (\codice{top}) \\
    Altezze & DSM assoluto e nDSM ufficiale & solo DSM assoluto (nDSM calcolato) \\
    Etichette & 6 classi, piene ed \emph{eroded} & 6 classi, piene ed \emph{eroded} \\
    Metrica principale & mIoU con tutte le classi & mIoU senza clutter (quasi assente) \\
    \hline
  \end{tabular}
\end{table}

\begin{table}[t]
  \centering
  \small
  \caption[Classi ISPRS e palette delle maschere]{Classi ISPRS, indici usati in questo lavoro e
    colori delle maschere.}
  \label{tab:classi}
  \begin{tabular}{|c|l|c|c|}
    \hline
    \textbf{Indice} & \textbf{Classe} & \textbf{RGB} & \textbf{Colore} \\
    \hline
    0 & Superfici impermeabili & (255, 255, 255) & \ClasseBoxX{ClsImpervious} \\
    1 & Edifici               & (0, 0, 255)     & \ClasseBoxX{ClsBuilding} \\
    2 & Vegetazione bassa     & (0, 255, 255)   & \ClasseBoxX{ClsLowVeg} \\
    3 & Alberi                & (0, 255, 0)     & \ClasseBoxX{ClsTree} \\
    4 & Auto                  & (255, 255, 0)   & \ClasseBoxX{ClsCar} \\
    5 & Clutter / sfondo      & (255, 0, 0)     & \ClasseBoxX{ClsClutter} \\
    255 & non etichettato (ignorato) & -- & -- \\
    \hline
  \end{tabular}
\end{table}

\subsection{Problemi noti dei file}

L'ispezione dei file distribuiti ha rivelato alcuni difetti che, se ignorati, falserebbero
silenziosamente i risultati (\cref{tab:difetti-file}). La preparazione dei dati li gestisce con
regole esplicite: un ordine di priorità tra le cartelle delle etichette, il rifiuto delle
etichette con più dello 0,5\% di colori fuori palette e un controllo delle dimensioni con
tolleranza di 2 pixel.

\begin{table}[t]
  \centering
  \small
  \caption[Difetti noti dei file ISPRS]{Difetti noti dei file ISPRS e loro gestione.}
  \label{tab:difetti-file}
  \begin{tabular}{|p{4.8cm}|p{2.9cm}|p{4.2cm}|}
    \hline
    \textbf{Difetto} & \textbf{Effetto se ignorato} & \textbf{Gestione} \\
    \hline
    Etichetta di \tile{4\_12} in \codice{5\_Labels\_all} corrotta (nessun colore ISPRS) &
      una tile di validation tutta ignorata & si usa la copia in \codice{5\_Labels\_for\_\allowbreak participants} \\
    \hline
    Etichetta di \tile{6\_7}: lo 0,7\% dei pixel \enquote{auto} ha colore (252, 255, 0) &
      pixel ignorati & stessa priorità tra cartelle \\
    \hline
    DSM e nDSM di \tile{3\_13} larghi 5999 pixel & disallineamento con l'immagine & una colonna replicata dal bordo \\
    \hline
    nDSM \codice{ownapproach}: 11 file su 38 vuoti & ablation impossibile & si usa la variante \codice{lastools} \\
    \hline
  \end{tabular}
\end{table}

% ---------------------------------------------------------------------
\section{Protocolli di split}
\label{sec:split}

Gli split sono sempre \textbf{per tile intere}: le tile vengono assegnate agli split \emph{prima}
del taglio in ritagli, quindi nessun ritaglio può contenere pixel di due split e nessuna porzione
di scena compare sia in training sia in valutazione~\cite{kaufman2012leakage}
(\cref{fig:split}).

\begin{description}[style=nextline, leftmargin=1.2em]
  \item[Potsdam \enquote{standard}.] È lo split di riferimento della
  letteratura~\cite{mmseg2020}: 24 tile di training, che coincidono con quelle distribuite con le
  etichette ai partecipanti del contest, e 14 tile di test. La validation usa due tile di
  training (\tile{4\_12} e \tile{7\_11}) tenute fuori per intero. Su questo protocollo si basano
  tutte le conclusioni.
  \item[Potsdam \enquote{tesi 2022}.] Riproduce il protocollo del 2022: esclude le stesse dieci
  aree, usa 22 tile di training e 6 di validation. Il testo originale non specifica quali sei:
  qui sono fissate a \tile{3\_14}, \tile{4\_11}, \tile{4\_15}, \tile{5\_13}, \tile{6\_15} e
  \tile{7\_11}. Il \enquote{test} coincide con la validation, valutata su tile intere. Serve solo
  al confronto con la \cref{tab:risultati-2022}.
  \item[Vaihingen \enquote{standard}.] 16 aree di training e 17 di test. La validation usa le
  aree 11 e 30. L'area 26 resta in training perché contiene gran parte del poco clutter
  disponibile.
\end{description}

\begin{figure}[t]
  \centering
  \resizebox{\linewidth}{!}{% Mappa delle 38 tile di Potsdam (righe 2-7, colonne 7-15) con l'assegnazione agli split:
% (a) protocollo standard della letteratura, (b) protocollo della tesi del 2022.
\begin{tikzpicture}[x=7.2mm, y=7.2mm]
  \newcommand{\cella}[4]{% riga, colonna, colore di riempimento, colore del testo
    \draw[black!45, line width=0.3pt, fill=#3] (#2-0.5,-#1-0.5) rectangle (#2+0.5,-#1+0.5);
    \node[font=\tiny, text=#4] at (#2,-#1) {#1\_#2};}
  \newcommand{\esclusa}[2]{%
    \draw[black!45, line width=0.3pt, pattern=north east lines, pattern color=black!35]
      (#2-0.5,-#1-0.5) rectangle (#2+0.5,-#1+0.5);
    \node[font=\tiny, text=black!60, fill=white, inner sep=0.5pt] at (#2,-#1) {#1\_#2};}
  \newcommand{\assi}{%
    \foreach \c in {7,...,15} \node[font=\tiny, text=SapGrigio] at (\c,-1.2) {\c};
    \foreach \r in {2,...,7} \node[font=\tiny, text=SapGrigio] at (6.0,-\r) {\r};}

  % ---------------- (a) standard ------------------------------------------------
  \begin{scope}
    \assi
    \foreach \r/\c in {2/10,2/11,2/12,3/10,3/11,3/12,4/10,4/11,5/10,5/11,5/12,
                       6/7,6/8,6/9,6/10,6/11,6/12,7/7,7/8,7/9,7/10,7/12}
      {\cella{\r}{\c}{SplitTrain}{white}}
    \foreach \r/\c in {4/12,7/11} {\cella{\r}{\c}{SplitVal}{black}}
    \foreach \r/\c in {2/13,2/14,3/13,3/14,4/13,4/14,4/15,5/13,5/14,5/15,6/13,6/14,6/15,7/13}
      {\cella{\r}{\c}{SplitTest}{white}}
    \node[font=\small\bfseries, anchor=west] at (5.6,-0.3) {(a) standard};
  \end{scope}

  % ---------------- (b) tesi del 2022 ------------------------------------------
  \begin{scope}[xshift=76mm]
    \assi
    \foreach \r/\c in {2/10,2/11,2/12,2/13,2/14,3/10,6/7,6/13,7/13,7/9} {\esclusa{\r}{\c}}
    \foreach \r/\c in {3/11,3/12,3/13,4/10,4/12,4/13,4/14,5/10,5/11,5/12,5/14,5/15,
                       6/8,6/9,6/10,6/11,6/12,6/14,7/7,7/8,7/10,7/12}
      {\cella{\r}{\c}{SplitTrain}{white}}
    \foreach \r/\c in {3/14,4/15,5/13,6/15,7/11,4/11} {\cella{\r}{\c}{SplitVal}{black}}
    \node[font=\small\bfseries, anchor=west] at (5.6,-0.3) {(b) tesi del 2022};
  \end{scope}

  % ---------------- legenda --------------------------------------------------
  \begin{scope}[shift={(6.5,-8.4)}, font=\footnotesize]
    \draw[black!45, fill=SplitTrain] (0,-0.25) rectangle (0.5,0.25); \node[anchor=west] at (0.6,0) {training};
    \draw[black!45, fill=SplitVal] (3.2,-0.25) rectangle (3.7,0.25); \node[anchor=west] at (3.8,0) {validation};
    \draw[black!45, fill=SplitTest] (6.8,-0.25) rectangle (7.3,0.25); \node[anchor=west] at (7.4,0) {test};
    \draw[black!45, pattern=north east lines, pattern color=black!35] (9.5,-0.25) rectangle (10.0,0.25);
    \node[anchor=west] at (10.1,0) {esclusa};
  \end{scope}
\end{tikzpicture}
}
  \caption[Split delle tile di Potsdam]{Assegnazione delle 38 tile di Potsdam agli split.
    (a)~Protocollo standard della letteratura, usato per tutte le conclusioni. (b)~Protocollo
    della tesi del 2022: quattro delle aree escluse (\tile{2\_13}, \tile{2\_14}, \tile{6\_13},
    \tile{7\_13}) sono tile di test nel protocollo standard.}
  \label{fig:split}
\end{figure}

% ---------------------------------------------------------------------
\section{Preparazione dei dati}
\label{sec:preparazione}

Le maschere RGB vengono convertite in indici con una tabella di conversione vettorizzata sui
codici a 24 bit dei colori. I colori fuori palette diventano 255. Le tile di training vengono
tagliate in ritagli di $512\times512$ pixel con passo 512. L'ultimo ritaglio di ogni riga e
colonna è allineato al bordo, e i ritagli completamente non etichettati vengono scartati. Su
Potsdam si ottengono 144 ritagli per tile, quindi \num{3168} ritagli di training (22 tile) e 288 di
validation (2 tile). Su Vaihingen, con sole 14 aree di training, i ritagli si sovrappongono
(passo 256), ma solo all'interno del training. Le tile di validation e di test vengono registrate
anche intere, con la maschera piena e quella eroded, per la valutazione finale.

\paragraph{Regime a poche etichette.}
Per la \RQ{3} il sottoinsieme di ritagli etichettati viene campionato \emph{per tile} con seed
fisso. A parità di epoche, una frazione piccola di etichette produrrebbe pochissimi
aggiornamenti dei pesi (con l'1\% delle etichette, 2 batch per epoca) e il confronto tra
pre-addestramenti misurerebbe soprattutto il sotto-addestramento. I ritagli etichettati vengono
quindi ripetuti, con un nuovo rimescolamento a ogni passata, fino ad avere per epoca tanti batch
quanti con tutte le etichette. Ogni frazione riceve così lo stesso budget di iterazioni
(\cref{tab:frazioni}). Il pre-addestramento DINO usa comunque tutte le immagini di training,
senza etichette.

\begin{table}[t]
  \centering
  \small
  \caption[Regime a poche etichette su Potsdam]{Regime a poche etichette su Potsdam: ritagli
    etichettati e batch per epoca (batch effettivo 8, accumulo 4), con e senza ripetizione fino
    alla lunghezza di un'epoca completa.}
  \label{tab:frazioni}
  \begin{tabular}{|c|r|c|c|}
    \hline
    \textbf{Frazione} & \makecell{\textbf{Ritagli}\\\textbf{etichettati}} & \makecell{\textbf{Batch/epoca}\\\textbf{(con ripetizione)}} & \makecell{\textbf{Batch/epoca}\\\textbf{(senza)}} \\
    \hline
    100\% & \num{3168} & 396 & 396 \\
    25\%  & 792  & 396 & 99 \\
    10\%  & 308  & 396 & 38 \\
    1\%   & 22   & 396 & 2 \\
    \hline
  \end{tabular}
\end{table}

\paragraph{Pool non etichettati.}
Per verificare se più dati di dominio aiutino DINO, il pre-addestramento può usare anche i frame
grezzi di Toronto (RGB a 8 bit, \num{4433} ritagli) e di Vaihingen (IR-R-G a 16 bit, riportati a
8 bit con uno stretch ai percentili 0,5--99,5). I frame di Vaihingen coprono anche le aree di
test di Vaihingen, quindi un pool della stessa città del dataset di valutazione viene rifiutato
automaticamente, per evitare un'informazione trasduttiva sul test.

% ---------------------------------------------------------------------
\section{Configurazione di addestramento}
\label{sec:addestramento}

\subsection{Segmentazione}

La \cref{tab:iperparametri} riporta le configurazioni di addestramento. La configurazione di base
mantiene gli iperparametri del 2022 dove possibile (learning rate, StepLR, batch effettivo 32,
100 epoche) e corregge i difetti del protocollo. L'ottimizzatore è AdamW~\cite{loshchilov2019adamw}
(nel 2022 Adam~\cite{kingma2015adam}). Il batch di 8 ritagli con accumulo del gradiente su 4 passi
dà un batch effettivo di 32 nei limiti di memoria della GPU, con precisione mista
bfloat16~\cite{micikevicius2018mixed}. La normalizzazione con media e deviazione standard di
ImageNet è il primo strato del modello: training e inferenza non possono divergere su questo
punto. L'augmentation di training applica a immagine e maschera, in modo identico, riflessioni
orizzontali e verticali e rotazioni di \ang{90}: sono gli elementi di $D_4$, che non richiedono
interpolazione e lasciano le etichette esatte. Il miglior checkpoint viene scelto sulla mIoU di
validation, e l'addestramento si ferma dopo 20 epoche senza miglioramenti.

\begin{table}[t]
  \centering
  \small
  \caption[Configurazioni di addestramento della segmentazione]{Configurazioni di addestramento
    della segmentazione. La colonna \enquote{tesi 2022} riproduce il protocollo originale ed è
    usata solo per il confronto con la \cref{tab:risultati-2022}.}
  \label{tab:iperparametri}
  \begin{tabular}{|l|c|c|c|}
    \hline
    \textbf{Parametro} & \textbf{Base (ResNet)} & \textbf{Transformer} & \textbf{Tesi 2022} \\
    \hline
    Ritaglio / passo & 512 / 512 & 512 / 512 & 256 / 256 \\
    Augmentation & $D_4$ & $D_4$ & nessuna \\
    Ottimizzatore & AdamW & AdamW & Adam \\
    Learning rate & $10^{-3}$ & $10^{-4}$ & $10^{-3}$ \\
    Weight decay & $10^{-4}$ & $10^{-4}$ & -- \\
    Scheduler & StepLR (15; 0,4) & coseno & StepLR (15; 0,4) \\
    Batch $\times$ accumulo & $8\times4$ & $4\times8$ & $32\times1$ \\
    Epoche / early stopping & 100 / 20 & 100 / 20 & 100 / no \\
    Precisione & bf16 mista & bf16 mista & bf16 mista \\
    Dropout / drop path & 0,1 / -- & 0,1 / 0,2 & 0,1 / -- \\
    Valutazione: finestra / passo & 512 / 256 & 512 / 256 & 256 / 128 \\
    \hline
  \end{tabular}
\end{table}

\paragraph{Quantizzazione.}
La \maresunetvq{} usa un codebook di $K=512$ codici con $\beta=0{,}25$, aggiornamento EMA con
$\gamma=0{,}99$, inizializzazione dai dati e riavvio dei codici con conteggio EMA inferiore a
0,1. Durante l'addestramento si registrano perplexity e uso del codebook a ogni epoca.
L'ablation della \RQ{2} varia $K\in\{64,256,512,1024\}$, $\beta\in\{0{,}1;\,0{,}25;\,1{,}0\}$,
il peso $\lambda$, la quantizzazione delle connessioni di salto e l'aggiornamento per gradiente
($\gamma=0$) al posto dell'EMA.

\paragraph{Encoder transformer.}
Oltre a ResNet-18/34/50, l'encoder può essere uno Swin-T~\cite{liu2021swin} o un
ViT~\cite{dosovitskiy2021vit}. Il ViT ha una sola scala, e le quattro scale richieste dal decoder
vengono costruite con una \emph{simple feature pyramid}~\cite{li2022vitdet}. Tutti i pesi
ImageNet vengono portati da torchvision, con l'equivalenza numerica verificata da test.

\subsection{Pre-addestramento DINO}

La \cref{tab:iperparametri-dino} riporta la configurazione del pre-addestramento e la
\cref{tab:augmentation-dino} le due politiche di augmentation confrontate nella \RQ{4}.

\begin{table}[t]
  \centering
  \small
  \caption[Configurazione del pre-addestramento DINO]{Configurazione del pre-addestramento DINO,
    a confronto con quanto dichiarato e quanto implementato nel 2022.}
  \label{tab:iperparametri-dino}
  \begin{tabular}{|l|c|c|}
    \hline
    \textbf{Parametro} & \textbf{Questo lavoro} & \textbf{2022 (codice)} \\
    \hline
    Testa di proiezione & \makecell{MLP 2048-2048-256\\$K=\num{16384}$ prototipi} & \makecell{assente\\(uscita densa)} \\
    Normalizzazione della softmax & per campione & sull'intero batch \\
    Viste & $2\times224^2+6\times96^2$ & viste locali a $224^2$ \\
    Learning rate & \makecell{$5\cdot10^{-4}\cdot B/256$\\warmup 10 epoche, coseno} & fisso \\
    Weight decay & coseno $0{,}04\to0{,}4$ & assente \\
    $\tau_s$ / $\tau_t$ & \makecell{0,1 / $0{,}04\to0{,}07$\\in 30 epoche} & 0,1 / 0,04 fissa \\
    Momentum del teacher & coseno $0{,}996\to1$ & 0,996 fisso \\
    Teacher iniziale & copia dello student & rete indipendente \\
    Centering & $\lambda=0{,}9$, per campione & \makecell{$\lambda=0{,}9$, sul batch\\appiattito} \\
    Epoche / batch & 100 / 32 & 100 / -- \\
    Clipping del gradiente & norma 3,0 & -- \\
    \hline
  \end{tabular}
\end{table}

\begin{table}[t]
  \centering
  \small
  \caption[Politiche di augmentation per DINO]{Politiche di augmentation per DINO confrontate
    nella \RQ{4}. Le probabilità di blur sono 1,0 e 0,1 per le due viste globali e 0,5 per le
    locali; la solarizzazione si applica solo alla seconda vista globale.}
  \label{tab:augmentation-dino}
  \begin{tabular}{|l|c|c|}
    \hline
    \textbf{Trasformazione} & \textbf{Standard} & \textbf{Aerea} \\
    \hline
    Riflessione orizzontale & sì & sì \\
    Riflessione verticale, rotazioni \ang{90} & no & sì \\
    Jitter ($p=0{,}8$): luminosità, contrasto & 0,4; 0,4 & 0,2; 0,2 \\
    Jitter: saturazione, tinta & 0,2; 0,1 & 0,1; 0,02 \\
    Scala di grigi & $p=0{,}2$ & $p=0{,}1$ \\
    Solarizzazione (2\textsuperscript{a} vista globale) & $p=0{,}2$ & mai \\
    Blur gaussiano & sì & sì \\
    \hline
  \end{tabular}
\end{table}

% ---------------------------------------------------------------------
\section{Protocollo di valutazione}
\label{sec:valutazione}

La valutazione finale avviene sulle tile di test intere, con la finestra scorrevole della
\cref{sec:finestra} (finestra 512, passo 256) e i pesi del miglior checkpoint di validation. Ogni
tile viene predetta una volta sola e confrontata con due riferimenti, le etichette piene e quelle
eroded. Si accumulano così due matrici di confusione sull'intero split. Dalle matrici si
ricavano, per ciascun riferimento, IoU, F1, precision e recall per classe, OA, mF1 e mIoU, con
tutte le classi e senza clutter. Il formato riprende le tabelle 4.3--4.8 della tesi del 2022, per
facilitare il confronto. La metrica principale è la mIoU con tutte le classi su Potsdam e la mIoU
senza clutter su Vaihingen.

% ---------------------------------------------------------------------
\section{Implementazione}
\label{sec:implementazione}

Il codice del 2022 non è stato riusato. Il pacchetto \rsseg{} (\emph{Remote Sensing
SEGmentation}) è una reimplementazione da zero (\cref{fig:pipeline}): modelli e addestramento in
Keras~3~\cite{chollet2015keras} su backend PyTorch~\cite{paszke2019pytorch}, tracciamento degli
esperimenti con MLflow~\cite{zaharia2018mlflow}, configurazioni YAML componibili con validazione
delle chiavi. I parametri, la configurazione completa, il commit del codice, le metriche per
epoca, il picco di memoria della GPU, i checkpoint e i report di ogni esecuzione vengono
registrati su MLflow.

\begin{figure}[t]
  \centering
  \resizebox{\linewidth}{!}{% Pipeline sperimentale: dai dati ISPRS al report, con il ramo di pre-addestramento DINO
% e l'orchestrazione delle campagne. Coordinate esplicite su una griglia 5 x 2.
\begin{tikzpicture}[
    p/.style={attn, text width=2.4cm, minimum height=14mm},
    ds/.style={dato, text width=2.3cm, minimum height=12mm, font=\scriptsize},
  ]
  \def\dx{3.25}
  % riga superiore: segmentazione supervisionata
  \node[ds] (raw)   at (0,0)      {Dati ISPRS\\TOP, etichette, DSM/nDSM};
  \node[p]  (prep)  at (\dx,0)    {\texttt{prepare}\\{\scriptsize split per tile, LUT, ritagli, nDSM, statistiche}};
  \node[ds] (man)   at (2*\dx,0)  {ritagli e tile intere\\\texttt{crops.csv}, \texttt{split.json}};
  \node[p]  (train) at (3*\dx,0)  {\texttt{train}\\{\scriptsize Keras \texttt{fit}, CE + VQ, early stopping}};
  \node[p]  (eval)  at (4*\dx,0)  {valutazione\\{\scriptsize finestra scorrevole, etichette piene + eroded}};
  % riga inferiore: pre-addestramento
  \node[ds] (pool)  at (0,-2.6)     {frame non etichettati\\Toronto, Vaihingen};
  \node[p]  (unl)   at (\dx,-2.6)   {\texttt{prepare-}\\\texttt{unlabeled}\\{\scriptsize ritagli 512, 8 bit}};
  \node[p]  (dino)  at (2*\dx,-2.6) {\texttt{pretrain}\\{\scriptsize DINO, multi-crop}};
  \node[ds] (bb)    at (3*\dx,-2.6) {encoder dello student\\\texttt{.weights.h5}};
  \node[blocco, fill=SapRossoVelo, text width=2.4cm, minimum height=12mm] (mlf) at (4*\dx,-2.6)
       {MLflow\\{\scriptsize parametri, metriche, artefatti}};
  % orchestrazione
  \node[quant, text width=8.6cm] (camp) at (2.5*\dx,2.0)
       {\texttt{campaign}: passi in sottoprocessi, prerequisiti automatici, ripresa, report};

  \draw[flusso] (raw) -- (prep);
  \draw[flusso] (prep) -- (man);
  \draw[flusso] (man) -- (train);
  \draw[flusso] (train) -- (eval);
  \draw[flusso] (pool) -- (unl);
  \draw[flusso] (unl) -- (dino);
  \draw[flusso] (man) -- node[nota, right, pos=0.45] {immagini di training} (dino);
  \draw[flusso] (dino) -- (bb);
  \draw[flusso] (bb) -- node[nota, right, pos=0.45] {inizializzazione} (train);
  \draw[flusso] (eval) -- (mlf);
  \draw[flusso, densely dotted] (bb) -- (mlf);

  \draw[->, SapRossoScuro, densely dashed] (camp.west) -| (prep.north);
  \draw[->, SapRossoScuro, densely dashed] (camp.south -| train.north) -- (train.north);
  \draw[->, SapRossoScuro, densely dashed] (camp.east) -| (eval.north);
\end{tikzpicture}
}
  \caption[Pipeline sperimentale]{Pipeline sperimentale. La preparazione divide i dati per tile
    e produce ritagli e tile intere. Il ramo inferiore pre-addestra l'encoder con DINO, anche con
    pool non etichettati. Le campagne orchestrano i passi in sottoprocessi riprendibili, e ogni
    passo registra parametri, metriche e artefatti su MLflow.}
  \label{fig:pipeline}
\end{figure}

\paragraph{Verifica.}
109 test automatici, eseguiti in integrazione continua, verificano le proprietà da cui dipendono
le conclusioni (\cref{tab:verifiche}). Alcune sono equivalenze numeriche con riferimenti
indipendenti, altre sono invarianti dei dati, altre ancora test di regressione sui difetti del
2022.

\begin{table}[t]
  \centering
  \small
  \caption[Verifiche automatiche]{Principali verifiche automatiche dell'implementazione.}
  \label{tab:verifiche}
  \begin{tabular}{|l|p{9.6cm}|}
    \hline
    \textbf{Ambito} & \textbf{Proprietà verificata} \\
    \hline
    Encoder & ResNet-18/34 Keras = torchvision (pesi e statistiche BN casuali), Swin e ViT con padding e shift \\
    Attenzione & PAM lineare = trascrizione della formula originale = attenzione esplicita $1+\hatq\T\hatk$; CAM = riferimento NumPy \\
    VQ & gradiente straight-through = identità; perdita solo in training; l'EMA ritrova centri noti con perplexity alta \\
    DINO & perdita = media delle perdite per campione; teacher = copia iniziale, non addestrabile, aggiornato per EMA \\
    Dati & split disgiunti; conversione colori biunivoca; allineamento immagine/maschera dopo l'augmentation; frazioni di etichette deterministiche \\
    nDSM & encoder esteso con pesi nulli = encoder RGB (\cref{prop:inflazione}) \\
    Metriche & metriche accumulate su più batch = calcolo NumPy sulla stessa matrice di confusione \\
    Regressione & il clutter riceve gradiente (\cref{prop:classe-esclusa}) \\
    End-to-end & dati sintetici $\to$ preparazione $\to$ training $\to$ report $\to$ DINO $\to$ fine-tuning \\
    Ripresa & un training messo in pausa e ripreso coincide con uno ininterrotto \\
    \hline
  \end{tabular}
\end{table}

\paragraph{Campagne.}
Ogni domanda di ricerca è una \emph{campagna} YAML che elenca configurazioni, seed, sweep di
parametri, pre-addestramenti richiesti e benchmark. L'esecutore espande la campagna in passi con
un identificativo stabile e prepara i prerequisiti mancanti (dati, pool, pre-addestramenti).
Esegue poi ogni passo in un sottoprocesso, così che un errore o un esaurimento della memoria
fermi solo quel passo. I passi completati vengono riusati tra campagne diverse, e una campagna
interrotta riprende dal punto in cui si era fermata. Training e pre-addestramenti possono essere
messi in pausa e ripresi in modo esatto: l'ordine dei dati e le augmentation dipendono solo dal
seed e dalla posizione nel flusso.

% ---------------------------------------------------------------------
\section{Budget computazionale}
\label{sec:budget}

Tutti gli esperimenti sono dimensionati per una singola GPU NVIDIA RTX 4060 da \SI{8}{\giga\byte}.
La \cref{tab:vram} riporta il picco di memoria e il throughput misurati su passi di training
reali. ResNet-50 con batch 8 non entra in memoria e viene eseguito con batch 4 e accumulo 8. Il
tempo di un'epoca su Potsdam è di circa 2,8 minuti per \maresunet{} R18 e 3,2 minuti per R34. La
valutazione delle 14 tile di test richiede circa 7 minuti. L'insieme di tutte le campagne è
stimato in circa 16 giorni di calcolo, di cui circa 4 per il solo ViT-B/16.

\begin{table}[t]
  \centering
  \small
  \caption[Memoria e throughput misurati su RTX 4060]{Picco di memoria della GPU e throughput
    misurati su RTX 4060 (\SI{8}{\giga\byte}) con passi di training reali, ritagli $512^2$ e
    precisione bf16 mista.}
  \label{tab:vram}
  \begin{tabular}{|l|l|r|r|}
    \hline
    \textbf{Configurazione} & \textbf{Batch $\times$ accumulo} & \textbf{VRAM} & \textbf{Immagini/s} \\
    \hline
    \maresunet{} R18          & $8\times4$ & \SI{4,8}{\giga\byte} & 21 \\
    \maresunet{} R18 + nDSM   & $8\times4$ & \SI{4,8}{\giga\byte} & 23 \\
    \maresunetvq{} R18        & $8\times4$ & \SI{4,7}{\giga\byte} & 23 \\
    \maresunet{} R34          & $8\times4$ & \SI{5,5}{\giga\byte} & 19 \\
    \maresunet{} R50          & $8\times4$ & esaurita & -- \\
    \maresunet{} Swin-T       & $4\times8$ & \SI{4,2}{\giga\byte} & 13 \\
    \maresunet{} ViT-B/16     & $4\times8$ & \SI{7,1}{\giga\byte} & 2 \\
    DINO R18 ($2\times224^2+6\times96^2$) & $32$ & \SI{4,2}{\giga\byte} & 85 \\
    DINO R34                  & $32$ & \SI{6,2}{\giga\byte} & 40 \\
    \hline
  \end{tabular}
\end{table}

% ---------------------------------------------------------------------
\section{Analisi statistica e riproducibilità}
\label{sec:statistica}

Le differenze attese tra le varianti sono dell'ordine di uno o due punti di mIoU, comparabili con
la variabilità tra esecuzioni con seed diversi~\cite{bouthillier2021variance}. Per questo:

\begin{itemize}
  \item ogni configurazione viene addestrata con \textbf{tre seed} (0, 1, 2), e si riportano media
  e deviazione standard;
  \item i confronti sono \textbf{appaiati per seed}: il seed fissa l'inizializzazione delle parti
  non pre-addestrate, l'ordine dei dati e le augmentation, quindi la differenza
  $\Delta^{(s)}=\mathrm{mIoU}^{(s)}_{\mathrm{variante}}-\mathrm{mIoU}^{(s)}_{\mathrm{rif}}$
  elimina una parte della variabilità comune;
  \item la variabilità dovuta al campione di test (solo 14 tile su Potsdam) viene stimata con un
  \textbf{bootstrap sulle tile}~\cite{efron1993bootstrap}: si ricampionano con reinserimento le
  matrici di confusione delle singole tile e si ricalcola la mIoU, ottenendo un intervallo al
  95\%;
  \item un effetto viene considerato sostenuto dai dati solo se $\Delta^{(s)}$ ha lo stesso segno
  per tutti i seed \emph{e} l'intervallo bootstrap della differenza esclude lo zero; altrimenti
  il risultato è riportato come non conclusivo.
\end{itemize}

\Segnaposto{Implementare il bootstrap per tile: il report salva oggi la mIoU di ogni tile ma
non la sua matrice di confusione, che serve per ricalcolare la mIoU globale sui campioni
bootstrap. Fissare il numero di ricampionamenti.}

Ogni esecuzione salva la configurazione completa, il commit del codice e il manifest dei dati, e
gli split sono deterministici e verificati da test. I risultati di ogni tabella del
\cref{cap:esperimenti} si possono quindi rigenerare con un solo comando di campagna.

% =====================================================================
%  Capitolo 5 — Esperimenti
%  Le tabelle hanno già righe, colonne e didascalie definitive: i valori \ND vanno
%  sostituiti con "media $\pm$ dev. std." sui tre seed, presi dai report delle campagne
%  (runs/campaigns/<campagna>/report.md) e da MLflow.
% =====================================================================
\chapter{Esperimenti}
\label{cap:esperimenti}

Questo capitolo raccoglie gli esperimenti, organizzati per domanda di ricerca secondo il disegno
della \cref{tab:disegno}. Salvo dove indicato, i valori sono mIoU, mF1 e OA sul test di Potsdam
(protocollo standard, etichette piene), come media $\pm$ deviazione standard su tre seed. Le
colonne \enquote{eroded} usano le etichette senza bordi.

\Segnaposto{Gli esperimenti completi non sono ancora stati eseguiti. Ogni sezione contiene le
tabelle e le figure previste, con valori \ND{} da sostituire. Non trarre conclusioni prima di
aver completato le campagne corrispondenti.}

% ---------------------------------------------------------------------
\section{Stato delle campagne}
\label{sec:stato-campagne}

\begin{table}[H]
  \centering
  \small
  \caption[Stato delle campagne di esperimenti]{Stato delle campagne di esperimenti. Le prove
    \enquote{rapide} verificano la catena completa su pochi passi e non producono risultati.}
  \label{tab:stato-campagne}
  \begin{tabular}{|l|c|l|l|}
    \hline
    \textbf{Campagna} & \textbf{Passi} & \textbf{Durata stimata} & \textbf{Stato} \\
    \hline
    \codice{thesis2022}      & 3   & $\sim$15 h            & da eseguire \\
    \codice{rq1\_attention}  & 6   & $\sim$18--29 h        & da eseguire \\
    \codice{thesis\_grid}    & 26  & $\sim$79--125 h       & da eseguire \\
    \codice{rq3\_low\_label} & 38  & $\sim$112--177 h      & da eseguire \\
    \codice{rq4\_aerial}     & 14  & $\sim$32--48 h        & da eseguire \\
    \codice{rq5\_efficiency} & 10  & $\sim$6 min           & prova rapida eseguita \\
    \codice{rq6\_ndsm}       & 9   & $\sim$27--44 h        & da eseguire \\
    \codice{transformers}    & 10  & $\sim$30--195 h       & da eseguire \\
    \codice{vaihingen}       & 9   & $\sim$8--13 h         & da eseguire \\
    \hline
  \end{tabular}
\end{table}

\Segnaposto{Aggiornare la colonna \enquote{Stato} (e le durate misurate) man mano che le
campagne vengono completate.}

% ---------------------------------------------------------------------
\section{Riproduzione del protocollo del 2022}
\label{sec:exp-2022}

Il primo esperimento addestra il modello migliore della tesi del 2022 (\maresunet{} R18 ImageNet)
sul protocollo di allora (\cref{fig:split}b), ma con le correzioni che non cambiano il protocollo:
clutter nella perdita, metriche globali, normalizzazione degli ingressi, rimescolamento dei
batch. Serve a separare l'effetto delle correzioni da quello del cambio di protocollo.

\begin{table}[H]
  \centering
  \small
  \caption[Confronto con la tesi del 2022 sul suo protocollo]{\maresunet{} R18 ImageNet sul
    protocollo della tesi del 2022: valori riportati allora e reimplementazione.}
  \label{tab:exp-2022}
  \begin{tabular}{|l|c|c|c|c|}
    \hline
    \textbf{Esecuzione} & \textbf{OA} & \textbf{mF1} & \textbf{mIoU} & \textbf{IoU clutter} \\
    \hline
    \rowcolor{RigaRiferimento}
    Tesi 2022 (tab.~4.5) & 0,791 & 0,722 & 0,604 & 0,149 \\
    Reimplementazione, protocollo 2022 & \ND & \ND & \ND & \ND \\
    Reimplementazione, protocollo standard & \ND & \ND & \ND & \ND \\
    \hline
  \end{tabular}
\end{table}

\Segnaposto{Inserire risultati esperimento \codice{thesis2022} (3 seed) e confrontarli con il
riferimento 2022. Verificare in particolare l'IoU del clutter e la distanza tra validation e test.}

% ---------------------------------------------------------------------
\section{\RQ{1}: contributo dell'attenzione lineare}
\label{sec:exp-rq1}

\begin{table}[H]
  \centering
  \small
  \caption[RQ1: connessioni di salto con e senza attenzione]{\RQ{1}. Encoder ResNet-18 ImageNet
    con connessioni di salto identità (rete di tipo U-Net) e con attenzione lineare
    (\maresunet{}).}
  \label{tab:exp-rq1}
  \begin{tabular}{|l|c|c|c|c|r|}
    \hline
    \textbf{Connessioni di salto} & \textbf{OA} & \textbf{mF1} & \textbf{mIoU} & \makecell{\textbf{mIoU}\\\textbf{eroded}} & \textbf{Parametri} \\
    \hline
    Identità (U-Net) & \ND & \ND & \ND & \ND & \ND \\
    LAM + CAM (\maresunet{}) & \ND & \ND & \ND & \ND & \ND \\
    \hline
    $\Delta$ appaiata (media $\pm$ dev. std.) & \multicolumn{2}{c|}{} & \ND & \ND & \\
    \hline
  \end{tabular}
\end{table}

\begin{table}[H]
  \centering
  \small
  \caption[RQ1: IoU per classe]{\RQ{1}. IoU per classe sul test di Potsdam.}
  \label{tab:exp-rq1-classi}
  \begin{tabular}{|l|c|c|c|c|c|c|}
    \hline
    \textbf{Modello} & \makecell{\ClasseBoxX{ClsImpervious}\\\textbf{Imp.}} & \makecell{\ClasseBoxX{ClsBuilding}\\\textbf{Edif.}} &
    \makecell{\ClasseBoxX{ClsLowVeg}\\\textbf{Veg. b.}} & \makecell{\ClasseBoxX{ClsTree}\\\textbf{Alberi}} &
    \makecell{\ClasseBoxX{ClsCar}\\\textbf{Auto}} & \makecell{\ClasseBoxX{ClsClutter}\\\textbf{Clutter}} \\
    \hline
    U-Net R18 & \ND & \ND & \ND & \ND & \ND & \ND \\
    \maresunet{} R18 & \ND & \ND & \ND & \ND & \ND & \ND \\
    \hline
  \end{tabular}
\end{table}

\Segnaposto{Inserire risultati esperimento \codice{rq1\_attention}. Analizzare le classi in cui la
differenza è maggiore alla luce della \cref{oss:contrasto} (coerenza delle grandi superfici
rispetto ai contorni) e misurare $\norm{\bar{\vect{m}}}_2$ sui modelli addestrati.}

% ---------------------------------------------------------------------
\section{\RQ{2}: quantizzazione vettoriale}
\label{sec:exp-rq2}

\subsection{Griglia della tesi sul protocollo standard}

\begin{table}[H]
  \centering
  \small
  \caption[RQ2: griglia architettura, encoder e pre-addestramento]{\RQ{2} (e \RQ{3} con tutte le
    etichette). Griglia della tesi del 2022 rieseguita sul protocollo standard.}
  \label{tab:exp-griglia}
  \begin{tabular}{|l|l|l|c|c|c|c|}
    \hline
    \textbf{Modello} & \textbf{Encoder} & \textbf{Pre-addestr.} & \textbf{OA} & \textbf{mF1} & \textbf{mIoU} & \textbf{PPL finale} \\
    \hline
    \maresunet{}   & R18 & ImageNet & \ND & \ND & \ND & -- \\
    \maresunet{}   & R34 & ImageNet & \ND & \ND & \ND & -- \\
    \maresunet{}   & R18 & DINO     & \ND & \ND & \ND & -- \\
    \maresunet{}   & R34 & DINO     & \ND & \ND & \ND & -- \\
    \maresunetvq{} & R18 & ImageNet & \ND & \ND & \ND & \ND \\
    \maresunetvq{} & R34 & ImageNet & \ND & \ND & \ND & \ND \\
    \maresunetvq{} & R18 & DINO     & \ND & \ND & \ND & \ND \\
    \maresunetvq{} & R34 & DINO     & \ND & \ND & \ND & \ND \\
    \hline
  \end{tabular}
\end{table}

\Segnaposto{Inserire risultati esperimento \codice{thesis\_grid} (26 passi).}

\subsection{Ablation del codebook}

\begin{table}[H]
  \centering
  \small
  \caption[RQ2: ablation del codebook]{\RQ{2}. Ablation del codebook su \maresunetvq{} R18
    ImageNet: dimensione $K$, peso del commitment $\beta$, aggiornamento e posizione della
    quantizzazione. La perplexity è misurata nell'ultima epoca.}
  \label{tab:exp-vq-ablation}
  \begin{tabular}{|c|c|l|l|c|c|c|}
    \hline
    $\bm{K}$ & $\bm{\beta}$ & \textbf{Aggiornamento} & \textbf{Quantizzato} & \textbf{mIoU} & \textbf{PPL} & \textbf{Uso} \\
    \hline
    64   & 0,25 & EMA & neck & \ND & \ND & \ND \\
    256  & 0,25 & EMA & neck & \ND & \ND & \ND \\
    \rowcolor{RigaRiferimento}
    512  & 0,25 & EMA & neck & \ND & \ND & \ND \\
    1024 & 0,25 & EMA & neck & \ND & \ND & \ND \\
    512  & 0,1  & EMA & neck & \ND & \ND & \ND \\
    512  & 1,0  & EMA & neck & \ND & \ND & \ND \\
    512  & 0,25 & gradiente ($\gamma=0$) & neck & \ND & \ND & \ND \\
    512  & 0,25 & EMA & neck + salto & \ND & \ND & \ND \\
    \hline
  \end{tabular}
\end{table}

\begin{figure}[H]
  \centering
  \FiguraMancante[4.5cm]{Perplexity e uso del codebook per epoca, con quantizzazione solo del
    neck e anche delle connessioni di salto (verifica della \cref{prop:aggiramento}).}
  \caption[RQ2: perplexity del codebook durante l'addestramento]{\RQ{2}. Perplexity e uso del
    codebook durante l'addestramento.}
  \label{fig:exp-perplexity}
\end{figure}

\Segnaposto{Inserire risultati esperimento di ablation VQ (sweep su \codice{model.vq.*}). Verificare
se la perplexity del neck diminuisce quando le connessioni di salto non sono quantizzate.}

% ---------------------------------------------------------------------
\section{\RQ{3}: pre-addestramento DINO con poche etichette}
\label{sec:exp-rq3}

\subsection{Diagnostica del pre-addestramento}

\begin{figure}[H]
  \centering
  \FiguraMancante[4.5cm]{Perdita DINO ed entropia media del teacher per epoca (R18 su Potsdam e
    su Potsdam + Toronto); riferimenti a $0$ e a $\log K\approx9{,}70$ per riconoscere il
    collasso (\cref{sec:collasso}).}
  \caption[RQ3: diagnostica del pre-addestramento DINO]{\RQ{3}. Diagnostica del
    pre-addestramento DINO.}
  \label{fig:exp-dino-diagnostica}
\end{figure}

\subsection{Frazione di etichette}

\begin{table}[H]
  \centering
  \small
  \caption[RQ3: mIoU al variare della frazione di etichette]{\RQ{3}. mIoU di test di
    \maresunet{} R18 al variare della frazione di etichette e del pre-addestramento
    dell'encoder. Stesso budget di iterazioni per ogni frazione (\cref{tab:frazioni}).}
  \label{tab:exp-rq3}
  \begin{tabular}{|c|c|c|c|}
    \hline
    \textbf{Etichette} & \textbf{ImageNet} & \makecell{\textbf{DINO}\\\textbf{(Potsdam)}} & \makecell{\textbf{DINO}\\\textbf{(Potsdam + Toronto)}} \\
    \hline
    1\%   & \ND & \ND & \ND \\
    10\%  & \ND & \ND & \ND \\
    25\%  & \ND & \ND & \ND \\
    100\% & \ND & \ND & \ND \\
    \hline
  \end{tabular}
\end{table}

\begin{figure}[H]
  \centering
  \FiguraMancante[5cm]{mIoU in funzione della frazione di etichette (asse logaritmico) per i
    tre pre-addestramenti, con barre di errore sui seed.}
  \caption[RQ3: curva di apprendimento rispetto alle etichette]{\RQ{3}. mIoU in funzione della
    frazione di etichette.}
  \label{fig:exp-rq3}
\end{figure}

\Segnaposto{Inserire risultati esperimento \codice{rq3\_low\_label} (38 passi). Valutare
l'\cref{ip:dino} con i criteri della \cref{sec:statistica}.}

% ---------------------------------------------------------------------
\section{\RQ{4}: augmentation per immagini aeree}
\label{sec:exp-rq4}

\begin{table}[H]
  \centering
  \small
  \caption[RQ4: augmentation standard e aerea]{\RQ{4}. Fine-tuning di \maresunet{} R18 da encoder
    DINO pre-addestrati con le due politiche della \cref{tab:augmentation-dino}.}
  \label{tab:exp-rq4}
  \begin{tabular}{|l|c|c|c|}
    \hline
    \textbf{Augmentation DINO} & \makecell{\textbf{mIoU}\\\textbf{(100\%)}} & \makecell{\textbf{mIoU}\\\textbf{(10\%)}} & \makecell{\textbf{Entropia finale}\\\textbf{del teacher}} \\
    \hline
    Standard & \ND & \ND & \ND \\
    Aerea    & \ND & \ND & \ND \\
    \hline
  \end{tabular}
\end{table}

\Segnaposto{Inserire risultati esperimento \codice{rq4\_aerial} (14 passi).}

% ---------------------------------------------------------------------
\section{\RQ{5}: accuratezza e costo}
\label{sec:exp-rq5}

\begin{table}[H]
  \centering
  \small
  \caption[RQ5: accuratezza e costo degli encoder]{\RQ{5}. Accuratezza e costo degli encoder.
    Memoria e throughput si riferiscono a un passo di training su RTX 4060 (\cref{tab:vram}); il
    tempo di inferenza è quello di una tile di test $6000\times6000$.}
  \label{tab:exp-rq5}
  \begin{tabular}{|l|r|r|r|r|c|}
    \hline
    \textbf{Encoder} & \textbf{Parametri} & \makecell{\textbf{GFLOP}\\\textbf{($512^2$)}} & \textbf{VRAM} & \textbf{s/tile} & \textbf{mIoU} \\
    \hline
    ResNet-18 & \ND & \ND & \SI{4,8}{\giga\byte} & \ND & \ND \\
    ResNet-34 & \ND & \ND & \SI{5,5}{\giga\byte} & \ND & \ND \\
    ResNet-50 (batch 4) & \ND & \ND & \ND & \ND & \ND \\
    Swin-T    & \ND & \ND & \SI{4,2}{\giga\byte} & \ND & \ND \\
    ViT-B/16  & \ND & \ND & \SI{7,1}{\giga\byte} & \ND & \ND \\
    ViT-S/16 (DINO) & \ND & \ND & \ND & \ND & \ND \\
    \hline
  \end{tabular}
\end{table}

\begin{figure}[H]
  \centering
  \FiguraMancante[5cm]{Diagramma accuratezza--costo (mIoU rispetto a tempo di inferenza per tile
    o GFLOP), con la frontiera di Pareto.}
  \caption[RQ5: frontiera accuratezza--costo]{\RQ{5}. Frontiera accuratezza--costo.}
  \label{fig:exp-rq5}
\end{figure}

\Segnaposto{Inserire risultati esperimenti \codice{rq5\_efficiency} (campagna completa) e
\codice{transformers}. Aggiungere il conteggio di parametri e FLOP e il tempo di inferenza per
tile.}

% ---------------------------------------------------------------------
\section{\RQ{6}: altezza sul terreno come quarto canale}
\label{sec:exp-rq6}

\begin{table}[H]
  \centering
  \small
  \caption[RQ6: RGB e RGB + nDSM]{\RQ{6}. \maresunet{} R18 ImageNet con ingresso RGB e RGB +
    nDSM (ufficiale \codice{lastools} e calcolato dal DSM).}
  \label{tab:exp-rq6}
  \begin{tabular}{|l|c|c|c|c|c|}
    \hline
    \textbf{Ingresso} & \textbf{mIoU} & \makecell{\textbf{IoU}\\\textbf{edif.}} & \makecell{\textbf{IoU}\\\textbf{imp.}} & \makecell{\textbf{IoU}\\\textbf{alberi}} & \makecell{\textbf{IoU}\\\textbf{veg. b.}} \\
    \hline
    RGB & \ND & \ND & \ND & \ND & \ND \\
    RGB + nDSM ufficiale & \ND & \ND & \ND & \ND & \ND \\
    RGB + nDSM calcolato & \ND & \ND & \ND & \ND & \ND \\
    \hline
  \end{tabular}
\end{table}

\begin{figure}[H]
  \centering
  \FiguraMancante[5cm]{Differenza tra le matrici di confusione normalizzate per riga (RGB + nDSM
    meno RGB), con evidenza delle coppie edifici/superfici impermeabili e alberi/vegetazione
    bassa.}
  \caption[RQ6: effetto dell'nDSM sulle confusioni tra classi]{\RQ{6}. Effetto dell'nDSM sulle
    confusioni tra classi.}
  \label{fig:exp-rq6}
\end{figure}

\Segnaposto{Inserire risultati esperimento \codice{rq6\_ndsm} e confronto tra inizializzazioni
\codice{mean} e \codice{zeros} del quarto canale (\cref{prop:inflazione}).}

% ---------------------------------------------------------------------
\section{Secondo benchmark: Vaihingen}
\label{sec:exp-vaihingen}

\begin{table}[H]
  \centering
  \small
  \caption[Vaihingen]{Risultati sul test di Vaihingen (17 aree). Metrica principale: mIoU senza
    clutter.}
  \label{tab:exp-vaihingen}
  \begin{tabular}{|l|c|c|c|}
    \hline
    \textbf{Modello} & \textbf{OA} & \makecell{\textbf{mF1}\\\textbf{(senza clutter)}} & \makecell{\textbf{mIoU}\\\textbf{(senza clutter)}} \\
    \hline
    U-Net R18 & \ND & \ND & \ND \\
    \maresunet{} R18 & \ND & \ND & \ND \\
    \maresunet{} R18 + nDSM calcolato & \ND & \ND & \ND \\
    \hline
  \end{tabular}
\end{table}

\Segnaposto{Inserire risultati esperimento \codice{vaihingen}.}

% ---------------------------------------------------------------------
\section{Analisi qualitativa}
\label{sec:exp-qualitativa}

\begin{figure}[H]
  \centering
  \FiguraMancante[6cm]{Pannelli per 3--4 porzioni di tile di test: ortofoto, verità a terra,
    predizione di \maresunet{} R18 ImageNet, predizione della variante migliore, mappa degli
    errori. Legenda con i colori ISPRS:
    \ClasseBoxX{ClsImpervious}~imp., \ClasseBoxX{ClsBuilding}~edif.,
    \ClasseBoxX{ClsLowVeg}~veg. bassa, \ClasseBoxX{ClsTree}~alberi,
    \ClasseBoxX{ClsCar}~auto, \ClasseBoxX{ClsClutter}~clutter.}
  \caption[Esempi qualitativi]{Esempi qualitativi sul test di Potsdam.}
  \label{fig:exp-qualitativa}
\end{figure}

\Segnaposto{Generare le figure con \codice{rsseg predict} e la visualizzazione dei pannelli,
scegliendo le porzioni prima di guardare le predizioni, per evitare una selezione favorevole.}

% ---------------------------------------------------------------------
\section{Sintesi per domanda di ricerca}
\label{sec:exp-sintesi}

\begin{table}[H]
  \centering
  \small
  \caption[Sintesi dei risultati per domanda di ricerca]{Sintesi dei risultati per domanda di
    ricerca, secondo i criteri della \cref{sec:statistica}: \enquote{sostenuta},
    \enquote{respinta} o \enquote{non conclusiva}.}
  \label{tab:exp-sintesi}
  \begin{tabular}{|c|l|c|c|}
    \hline
    \textbf{RQ} & \textbf{Ipotesi} & \textbf{$\Delta$ mIoU} & \textbf{Esito} \\
    \hline
    \RQ{1} & \Cref{ip:lam}   & \ND & \ND \\
    \RQ{2} & \Cref{ip:vq}    & \ND & \ND \\
    \RQ{3} & \Cref{ip:dino}  & \ND & \ND \\
    \RQ{4} & \Cref{ip:aug}   & \ND & \ND \\
    \RQ{5} & descrittiva     & --  & \ND \\
    \RQ{6} & \Cref{ip:ndsm}  & \ND & \ND \\
    \hline
  \end{tabular}
\end{table}

\Segnaposto{Compilare la sintesi solo a campagne concluse. Ogni esito va motivato con i valori
appaiati per seed e con l'intervallo bootstrap.}

% =====================================================================
%  Capitolo 6 — Prospettive e struttura finale attesa
%  Nessuna conclusione: gli esperimenti non sono ancora completi.
% =====================================================================
\chapter{Prospettive e lavoro futuro}
\label{cap:prospettive}

Gli esperimenti di questa tesi non sono ancora completi, quindi il capitolo \emph{non} trae
conclusioni sulle domande di ricerca. Riassume invece lo stato del lavoro, elenca le attività
rimanenti, discute le direzioni di ricerca aperte dal quadro teorico e fissa la struttura attesa
del capitolo conclusivo.

\section{Stato del lavoro}
\label{sec:stato-lavoro}

\begin{table}[H]
  \centering
  \small
  \caption[Stato del lavoro]{Stato delle componenti del lavoro.}
  \label{tab:stato-lavoro}
  \begin{tabular}{|p{5.6cm}|l|}
    \hline
    \textbf{Componente} & \textbf{Stato} \\
    \hline
    Lettura critica della tesi del 2022 & completata (\cref{cap:premesse}) \\
    Quadro teorico & completato (\cref{cap:teoria}) \\
    Protocollo sperimentale & definito (\cref{cap:metodologia}) \\
    Implementazione e test automatici & completati, 109 test \\
    Preparazione di Potsdam, Vaihingen e pool di Toronto & completata \\
    Benchmark di memoria su RTX 4060 & completato (\cref{tab:vram}) \\
    Bootstrap per tile nel report & da implementare \\
    Campagne di esperimenti & da eseguire (\cref{tab:stato-campagne}) \\
    Analisi dei risultati e conclusioni & da scrivere \\
    \hline
  \end{tabular}
\end{table}

\section{Attività rimanenti}
\label{sec:attivita}

L'ordine di esecuzione segue la dipendenza tra le domande e il costo delle campagne:

\begin{enumerate}
  \item \codice{thesis2022} e \codice{rq1\_attention}: fissano il riferimento e quantificano
  l'effetto delle sole correzioni di protocollo (\cref{sec:exp-2022,sec:exp-rq1});
  \item \codice{thesis\_grid}: pre-addestramenti DINO R18/R34 e griglia della tesi sul protocollo
  standard, che alimenta \RQ{2} e il punto al 100\% della \RQ{3};
  \item \codice{rq6\_ndsm} e \codice{vaihingen}: relativamente brevi, verificano la
  generalizzazione su un secondo benchmark;
  \item \codice{rq3\_low\_label} e \codice{rq4\_aerial}: le campagne più lunghe, che riusano i
  pre-addestramenti già eseguiti;
  \item \codice{rq5\_efficiency} completa e \codice{transformers}, valutando se escludere
  ViT-B/16 (circa 30 minuti per epoca sulla GPU disponibile);
  \item analisi statistica (seed appaiati e bootstrap), figure qualitative, sintesi.
\end{enumerate}

\Segnaposto{Aggiornare l'elenco con le date effettive di esecuzione e con eventuali campagne
aggiunte o tolte.}

\section{Prospettive di ricerca}
\label{sec:prospettive-ricerca}

Indipendentemente dall'esito delle campagne, il quadro teorico del \cref{cap:teoria} indica
alcune direzioni che questa tesi non esplora fino in fondo.

\paragraph{Auto-supervisione densa.}
L'\cref{oss:globale-denso} osserva che la perdita di DINO agisce dopo una media globale, mentre
la segmentazione è un compito denso. Una variante con perdita locale, che confronti le
caratteristiche di regioni corrispondenti tra viste diverse, allineerebbe l'obiettivo di
pre-addestramento al compito a valle. Lo stesso vale per la ricostruzione
mascherata~\cite{he2022mae,cong2022satmae}, che opera per costruzione sulle singole porzioni
dell'immagine.

\paragraph{Diagnostica degli encoder pre-addestrati.}
Una valutazione lineare o $k$-NN sull'encoder congelato separerebbe la qualità della
rappresentazione dall'effetto del fine-tuning. L'informazione mutua~\eqref{eq:dino-mi}, stimata a
ogni epoca, darebbe un criterio di arresto più informativo dell'entropia del solo teacher.

\paragraph{Modelli di fondazione.}
Encoder pre-addestrati su collezioni molto grandi e diverse, come DINOv2~\cite{oquab2024dinov2} o
modelli specifici per il telerilevamento~\cite{wang2022ssl_rs}, sono un riferimento naturale per
la \RQ{3}: dicono se il limite è nella quantità di dati di dominio o nell'obiettivo di
pre-addestramento.

\paragraph{Quantizzazione oltre il collo di bottiglia.}
Se la \cref{prop:aggiramento} trova conferma empirica, la quantizzazione delle connessioni di
salto, o una quantizzazione gerarchica come in VQ-VAE-2~\cite{razavi2019vqvae2}, è il modo per
verificare davvero l'ipotesi di un vantaggio delle rappresentazioni discrete. I codici appresi
potrebbero anche essere interpretati come \enquote{parole visive} del dominio aereo e confrontati
con le classi.

\paragraph{Attenzione efficiente.}
L'\cref{oss:contrasto} suggerisce che il LAM aggreghi contesto globale con poca selettività.
Nuclei di somiglianza più contrastati, ma ancora fattorizzabili
come~\eqref{eq:attenzione-lineare-generale}, o attenzioni locali a finestra come in
Swin~\cite{liu2021swin}, sono alternative con costo lineare da confrontare a parità di budget.

\paragraph{Fusione multimodale e generalizzazione.}
Oltre alla fusione precoce dell'nDSM, la fusione tardiva a due rami~\cite{audebert2018beyond} e il
trasferimento tra dataset (Potsdam IR-R-G $\to$ Vaihingen, stesso falso colore) misurerebbero
quanto i modelli generalizzano a città e sensori diversi.

\paragraph{Efficienza e inferenza a bordo.}
La motivazione energetica del 2022 resta valida. Quantizzazione dei pesi, distillazione verso
encoder più piccoli~\cite{hinton2015distilling} e misure di energia per tile, oltre che di
tempo, renderebbero la \RQ{5} direttamente utile per l'inferenza su piattaforme con risorse
limitate.

\section{Struttura attesa del capitolo conclusivo}
\label{sec:struttura-conclusioni}

A esperimenti completati, questo capitolo verrà sostituito da un capitolo di conclusioni con la
struttura seguente.

\begin{enumerate}
  \item \textbf{Sintesi del lavoro}: problema, revisione del 2022, protocollo e implementazione.
  \item \textbf{Risposte alle domande di ricerca}: per ciascuna, l'esito della
  \cref{tab:exp-sintesi}, con l'entità dell'effetto, la sua variabilità e le condizioni in cui
  vale.
  \item \textbf{Confronto con la tesi del 2022}: quali conclusioni di allora reggono con il
  protocollo corretto, quali cambiano e per quale motivo.
  \item \textbf{Limiti}: una sola GPU e tre seed, due soli dataset urbani europei, pre-addestramento
  limitato a DINO, valutazione di parte dei fattori solo su Potsdam.
  \item \textbf{Sviluppi futuri}: le direzioni della \cref{sec:prospettive-ricerca} ordinate in
  base ai risultati ottenuti.
\end{enumerate}

\Segnaposto{Scrivere le conclusioni solo dopo aver completato il \cref{cap:esperimenti}.}


\backmatter
\cleardoublepage
\phantomsection
\addcontentsline{toc}{chapter}{\bibname}
\bibliographystyle{sapthesis-it}
\bibliography{bibliografia}

\end{document}
